\documentclass[11pt]{article}

\usepackage[T1]{fontenc}
\usepackage[utf8]{inputenc}
\usepackage[margin=1in]{geometry}
\usepackage{amsmath,amssymb,amsthm,mathtools}
\usepackage[colorlinks=true,linkcolor=blue,citecolor=blue,urlcolor=blue]{hyperref}

\numberwithin{equation}{section}
\newtheorem{theorem}{Theorem}[section]
\newtheorem{proposition}[theorem]{Proposition}
\newtheorem{lemma}[theorem]{Lemma}
\newtheorem{corollary}[theorem]{Corollary}
\newtheorem{conjecture}[theorem]{Conjecture}
\newtheorem{problem}[theorem]{Problem}
\theoremstyle{definition}
\newtheorem{definition}[theorem]{Definition}
\newtheorem{example}[theorem]{Example}
\theoremstyle{remark}
\newtheorem{remark}[theorem]{Remark}

\DeclareMathOperator{\Var}{Var}
\DeclareMathOperator{\Prob}{P}
\DeclareMathOperator{\supp}{supp}
\newcommand{\E}{\mathbb E}
\newcommand{\N}{\mathbb N}
\newcommand{\Z}{\mathbb Z}
\newcommand{\1}{\mathbf 1}

\title{Zeta-Law Entropy, Gamma Curvature, and Laplace-Measure Applications}
\author{Pudim AI Project}
\date{Version v009 -- 2026-05-31}

\begin{document}
\maketitle

\begin{abstract}
This staged paper presents a unified public theory built around Laplace-transform
positivity: zeta-law entropy is reorganized by logarithmic compression,
gamma-ratio and polygamma inequalities are expressed through completely monotone
kernels, and new applications are obtained by transporting measure positivity
through Stieltjes, complete-Bernstein, stable-subordination, and generalized-
Stieltjes operations.  Compared with version v008, this version adds ten fully
solved applications.  They include a Stieltjes representation for a Ramanujan
integral, a stable-subordination representation of a Prabhakar measure, a
derivative-level Mills-ratio bound family, two negative answers to over-strong
Bernstein and complete-monotonic-degree conjectures, exact threshold theorems
for tau, W-symbol, and digamma-cutoff problems, a proof of a Gurland
logarithmic-complete-monotonicity conjecture, and a triangular compression
formula for Sokal's generalized-Stieltjes expressions.  The application ledger
now records twenty-three public applications, all with uniform status
\emph{Solved}.
\end{abstract}

\noindent\textbf{Keywords.} Riemann zeta function; zeta probability law; entropy; Laplace transform; complete monotonicity; Bernstein functions; Stieltjes functions; gamma function; digamma function; Mills ratio; generalized Stieltjes functions.

\section{Motivation and main results}

Many classical identities for \(\zeta(s)\) become more transparent after a probabilistic normalization.  For \(\beta>1\), the weights \(n^{-\beta}\) define a Gibbs law whose partition function is precisely \(\zeta(\beta)\).  The energy is \(\log n\), the free energy is \(\log\zeta(\beta)\), and ratios such as \(\zeta(\beta+r)/\zeta(\beta)\) become moments.  This viewpoint is elementary, but it packages several unrelated-looking zeta inequalities into a common calculus.

The paper develops six connected layers.  The moment layer records basic zeta-law identities.  The modular entropy layer compares the successor map \(n\mapsto n+1\) with its finite residue shadows modulo \(q\).  The Euler-score layer rewrites zeta energy through von Mangoldt divisibility events.  The Mellin-Planck layer treats \(\Gamma(s)\zeta(s)\) as a continuous partition function and converts generalized Holder inequalities into integral norm inequalities.  The tail layer studies \(\zeta_n(s)^{-1}\) as a reciprocal partition problem for the truncated zeta law.  The v009 Laplace-transport layer then carries the same positivity language into Stieltjes, complete-Bernstein, stable-subordination, determinant, and generalized-Stieltjes applications.

Section~\ref{sec:applications} records twenty-three solved external problems with stable application labels.  APP-0001 through APP-0013 are the earlier zeta-law, Gamma, digamma, tail, beta-window, and q-log-convexity entries.  APP-0014 through APP-0023 form the new v009 block: a Ramanujan integral Stieltjes classification, a Prabhakar stable-subordination formula, a Mills-ratio derivative hierarchy, two counterexamples to over-strong Bernstein or degree claims, exact tau, W-symbol, and digamma cutoffs, a Gurland log-complete-monotonicity proof, and a Sokal triangular compression formula.  The added results are staged as one measure-transport story rather than as a list of unrelated inequalities.

\section{Applications}
\label{sec:applications}

The public application ledger contains twenty-three solved entries.  The first
thirteen were staged in earlier theory versions.  The v009 layer adds ten more
entries, all unconditional and all reducible to the same structural principle:
find the positive kernel, then transport it through an operation that preserves
complete monotonicity or one of its Stieltjes/Bernstein descendants.
\begin{enumerate}
\item \textbf{APP-0001.} Alzer--Kwong convexity and concavity problem.  Solved by Theorem~\ref{thm:alzer-kwong}; source problem reference \cite{alzer-kwong}.
\item \textbf{APP-0002.} Nantomah zeta positivity problem.  Solved by Theorem~\ref{thm:nantomah-positivity}; source problem reference \cite{nantomah}.
\item \textbf{APP-0003.} Sroysang generalized Holder problem.  Solved by Theorem~\ref{thm:sroysang-holder}; source problem reference \cite{sroysang}.
\item \textbf{APP-0004.} Complete monotonicity of \((\log\Gamma(x)+\log\Gamma(1/x))''\).  Solved by Theorem~\ref{thm:reciprocal-gamma-complete-monotonicity}; source problem reference \cite{qi-lim-nantomah-polygamma}.
\item \textbf{APP-0005.} Exact inverse-tail floor formula at \(s=7\).  Solved by Theorem~\ref{thm:tail-s7-floor}; source context \cite{kim-song-tail,pan-wu-tail}.
\item \textbf{APP-0006.} Exact inverse-tail floor formula at \(s=8\).  Solved by Theorem~\ref{thm:tail-s8-floor}; source context \cite{kim-song-tail,pan-wu-tail}.
\item \textbf{APP-0007.} Complete monotonicity of \(-(\psi(x)\psi(1/x))''\).  Solved by Theorem~\ref{thm:reciprocal-digamma-product-complete-monotonicity}; source problem reference \cite{qi-lim-nantomah-polygamma}.
\item \textbf{APP-0008.} Counterexample to complete monotonicity of higher-order polygamma product curvature.  Solved by Theorem~\ref{thm:pn-complete-monotonicity-counterexample}; source problem reference \cite{qi-lim-nantomah-polygamma}.
\item \textbf{APP-0009.} Sharp reciprocal Gamma-product monotonicity threshold.  Solved by Theorem~\ref{thm:gamma-product-sharp-threshold}; source problem reference \cite{bulboaca-zayed-gamma}.
\item \textbf{APP-0010.} Nielsen \(k\)-beta derivative-ratio monotonicity.  Solved by Theorem~\ref{thm:nielsen-k-beta-derivative-ratio}; source problem reference \cite{yin-zhang-k-beta}.
\item \textbf{APP-0011.} Exact \(n=2\) beta-window endpoint.  Solved by Theorem~\ref{thm:q2-exact-endpoint}; source problem reference \cite{qi-lim-nantomah-polygamma}.
\item \textbf{APP-0012.} Baricz \(V_q\) strict \(q\)-log-convexity.  Solved by Theorem~\ref{thm:baricz-vq-strict-logconvexity}; source problem reference \cite{baricz-turan-thesis}.
\item \textbf{APP-0013.} Yin \((p,k)\)-digamma sharp \(\alpha\)-necessity.  Solved by Theorem~\ref{thm:yin-pk-alpha-necessity}; source problem reference \cite{yin-pk-digamma}.
\item \textbf{APP-0014.} Ramanujan integral Stieltjes/complete-Bernstein classification.  Solved by Theorem~\ref{thm:app14-ramanujan-stieltjes}; source problem reference \cite{mishra-swaminathan-ramanujan}.
\item \textbf{APP-0015.} Prabhakar Q-measure stable-subordination representation.  Solved by Theorem~\ref{thm:app15-prabhakar-subordination}; source problem reference \cite{sibisi-prabhakar}.
\item \textbf{APP-0016.} Mills-ratio all-derivative-level bound extractor.  Solved by Theorem~\ref{thm:app16-mills-all-l}; source problem reference \cite{from-mills}.
\item \textbf{APP-0017.} Gamma-quotient Bernstein-for-all counterexample.  Solved by Theorem~\ref{thm:app17-baricz-counterexample}; source problem reference \cite{baricz-turan}.
\item \textbf{APP-0018.} Qi--Guo tau-threshold supremum correction.  Solved by Theorem~\ref{thm:app18-qg-tau}; source problem reference \cite{qi-guo-2004}.
\item \textbf{APP-0019.} Wakrim W-symbol exact Bernstein range.  Solved by Theorem~\ref{thm:app19-w-symbol}; source problem reference \cite{wakrim-w}.
\item \textbf{APP-0020.} Qi \(h_\lambda\) degree-four conjecture refutation.  Solved by Theorem~\ref{thm:app20-qi-hlambda}; source problem reference \cite{qi-hlambda}.
\item \textbf{APP-0021.} Szabo digamma-cutoff theorem.  Solved by Theorem~\ref{thm:app21-szabo-cutoff}; source problem reference \cite{szabo-cm}.
\item \textbf{APP-0022.} Chen--Choi Gurland gamma-ratio logarithmic complete monotonicity.  Solved by Theorem~\ref{thm:app22-gurland-logcm}; source problem reference \cite{chen-choi-gurland}.
\item \textbf{APP-0023.} Sokal generalized-Stieltjes lambda-derivative compression.  Solved by Theorem~\ref{thm:app23-sokal-lambda}; source problem reference \cite{sokal-gs}.
\end{enumerate}

\section{Notation and zeta-law calculus}

\begin{definition}[Riemann zeta probability law]
\label{def:zeta-law}
For \(\beta>1\), define a probability law \(\rho_\beta\) on \(\N\) by
\begin{equation}
\label{eq:zeta-law}
\rho_\beta(n)=\frac{n^{-\beta}}{\zeta(\beta)},\qquad n\in\N.
\end{equation}
Equivalently, with \(E(n)=\log n\) and \(Z(\beta)=\zeta(\beta)\),
\begin{equation}
\label{eq:gibbs-form}
\rho_\beta(n)=\frac{e^{-\beta E(n)}}{Z(\beta)}.
\end{equation}
\end{definition}

\begin{definition}[Zeta free energy]
\label{def:zeta-free-energy}
The zeta free energy is
\begin{equation}
\label{eq:zeta-free-energy}
A(\beta)=\log\zeta(\beta),\qquad \beta>1.
\end{equation}
\end{definition}

\begin{proposition}[Free-energy derivatives]
\label{prop:free-energy-derivatives}
For every \(\beta>1\),
\begin{align}
\label{eq:free-energy-derivatives}
A'(\beta)&=-\E_\beta[\log N],&
A''(\beta)&=\Var_\beta(\log N).
\end{align}
\end{proposition}

\begin{proof}
Absolute convergence of the differentiated Dirichlet series on compact subintervals of \((1,\infty)\) permits termwise differentiation.  Thus
\begin{equation}
\label{eq:first-free-derivative-proof}
A'(\beta)=\frac{\zeta'(\beta)}{\zeta(\beta)}
=-\frac{\sum_{n\ge1}(\log n)n^{-\beta}}{\zeta(\beta)}
=-\E_\beta[\log N].
\end{equation}
Differentiating once more gives
\begin{equation}
\label{eq:second-free-derivative-proof}
A''(\beta)=\E_\beta[(\log N)^2]-\E_\beta[\log N]^2
=\Var_\beta(\log N).
\end{equation}
\end{proof}

The same source asks whether \(P_0(x)=\psi(x)\psi(1/x)\) is concave, and more strongly whether \(-P_0''\) is completely monotonic on \((0,\infty)\).  The next theorem answers the stronger question.

\begin{theorem}[APP-0007: Complete monotonicity of reciprocal digamma product curvature]
\label{thm:reciprocal-digamma-product-complete-monotonicity}
Let
\begin{equation}
\label{eq:reciprocal-digamma-p0}
P_0(x)=\psi(x)\psi(1/x),\qquad x>0.
\end{equation}
Then \(-P_0''\) is strictly completely monotonic on \((0,\infty)\).  In particular, \(P_0\) is strictly concave on \((0,\infty)\).
\end{theorem}

\begin{proof}
Direct differentiation gives
\begin{equation}
\label{eq:p0-second-derivative}
\begin{aligned}
P_0''(x)
=&\psi''(x)\psi(1/x)
-\frac{2}{x^2}\psi'(x)\psi'(1/x)\\
&+\frac{2}{x^3}\psi(x)\psi'(1/x)
+\frac{1}{x^4}\psi(x)\psi''(1/x).
\end{aligned}
\end{equation}
Use the shifted Weierstrass expansion
\begin{equation}
\label{eq:digamma-shifted-expansion}
\psi(x)=-\gamma-\frac1x+\sum_{m=1}^{\infty}a_m(x),
\qquad
a_m(x)=\frac{x}{m(m+x)},
\end{equation}
and, after substituting \(1/x\),
\begin{equation}
\label{eq:reciprocal-digamma-shifted-expansion}
\psi(1/x)=-\gamma-x+\sum_{n=1}^{\infty}b_n(x),
\qquad
b_n(x)=\frac1{n(nx+1)}.
\end{equation}
On compact subintervals of \((0,\infty)\), the series for \(a_m\), \(b_n\), and their first two derivatives are uniformly summable; the differentiated double products are dominated by products of summable majorants.  Hence \(P_0\) may be differentiated term by term twice.

Let \(A_0(x)=-\gamma-1/x\) and \(B_0(x)=-\gamma-x\).  The base term satisfies
\begin{equation}
\label{eq:p0-base-term}
-\frac{d^2}{dx^2}\bigl(A_0(x)B_0(x)\bigr)=-\frac{2\gamma}{x^3}.
\end{equation}
For \(n\ge1\), elementary rational differentiation and Laplace inversion give
\begin{equation}
\label{eq:p0-a0bn-laplace}
-\frac{d^2}{dx^2}\bigl(A_0(x)b_n(x)\bigr)
=\int_0^\infty e^{-xt}t^2\frac{n+(\gamma-n)e^{-t/n}}{n^2}\,dt,
\end{equation}
and, for \(m\ge1\),
\begin{equation}
\label{eq:p0-b0am-laplace}
-\frac{d^2}{dx^2}\bigl(B_0(x)a_m(x)\bigr)
=\int_0^\infty e^{-xt}t^2(m-\gamma)e^{-mt}\,dt.
\end{equation}
For \(mn>1\),
\begin{equation}
\label{eq:p0-ambn-laplace}
-\frac{d^2}{dx^2}\bigl(a_m(x)b_n(x)\bigr)
=\int_0^\infty e^{-xt}t^2
\left[
\frac{e^{-t/n}}{mn^2(mn-1)}
-\frac{e^{-mt}}{n(mn-1)}
\right]dt.
\end{equation}
The exceptional product is \(a_1(x)b_1(x)=x/(x+1)^2\), and
\begin{equation}
\label{eq:p0-exceptional-laplace}
-\frac{d^2}{dx^2}\frac{x}{(x+1)^2}
=\int_0^\infty e^{-xt}t^2(t-1)e^{-t}\,dt.
\end{equation}
The identities above follow from
\begin{equation}
\label{eq:p0-laplace-powers}
\frac1{(x+c)^3}=\int_0^\infty e^{-xt}\frac{t^2}{2}e^{-ct}\,dt,
\qquad
\frac1{(x+c)^4}=\int_0^\infty e^{-xt}\frac{t^3}{6}e^{-ct}\,dt.
\end{equation}

The coefficients of the double products group by the two identities
\begin{equation}
\label{eq:p0-group-n}
\sum_{m=1}^{\infty}\frac1{m(mn-1)}
=-\psi(1-1/n)-\gamma
\qquad(n\ge2)
\end{equation}
and
\begin{equation}
\label{eq:p0-group-m}
\sum_{n=1}^{\infty}\frac1{n(mn-1)}
=-\psi(1-1/m)-\gamma
\qquad(m\ge2).
\end{equation}
Both are immediate from \(\psi(z+1)+\gamma=\sum_{k=1}^\infty z/(k(k+z))\), with \(z=-1/n\) and \(z=-1/m\).  Combining \eqref{eq:p0-base-term}--\eqref{eq:p0-group-m} gives
\begin{equation}
\label{eq:p0-laplace-kernel-representation}
-P_0''(x)=\int_0^\infty e^{-xt}K(t)\,dt,
\end{equation}
where
\begin{equation}
\label{eq:p0-kernel}
\begin{aligned}
K(t)=t^2\Bigg[
&1-\gamma+(t-1)e^{-t}\\
&+\sum_{n=2}^{\infty}
\frac{n(1-e^{-t/n})-\psi(1-1/n)e^{-t/n}}{n^2}
+\sum_{m=2}^{\infty}\bigl(m+\psi(1-1/m)\bigr)e^{-mt}
\Bigg].
\end{aligned}
\end{equation}

It remains to prove \(K(t)>0\).  Since \(\psi\) is increasing and \(1-1/n\in[1/2,1)\),
\begin{equation}
\label{eq:p0-minus-psi-positive}
-\psi(1-1/n)>0\qquad(n\ge2).
\end{equation}
Also
\begin{equation}
\label{eq:p0-m-sum-positive}
m+\psi(1-1/m)\ge 2-\gamma-2\log2>0
\qquad(m\ge2),
\end{equation}
because \(\psi(1-1/m)\ge\psi(1/2)=-\gamma-2\log2\).  Thus the two sums in \eqref{eq:p0-kernel} are positive term by term, apart from the possible need to dominate \(1-\gamma+(t-1)e^{-t}\).

For \(n\ge2\),
\begin{equation}
\label{eq:p0-minus-psi-lower}
-\psi(1-1/n)
=\gamma+\int_{1-1/n}^{1}\psi'(u)\,du
>\gamma+\int_{1-1/n}^{1}\frac{du}{u^2}
=\gamma+\frac1{n-1}.
\end{equation}
Consequently the \(n\)-sum in \eqref{eq:p0-kernel} is bounded below by
\begin{equation}
\label{eq:p0-n-sum-lower}
Ae^{-t/2},
\qquad
A=\gamma(\zeta(2)-1)+2-\zeta(2).
\end{equation}
Therefore \(K(t)/t^2\) is bounded below by
\begin{equation}
\label{eq:p0-f-lower}
f(t)=1-\gamma+(t-1)e^{-t}+Ae^{-t/2}.
\end{equation}
Now
\begin{equation}
\label{eq:p0-f-endpoints}
f(0)=A-\gamma=(2-\zeta(2))(1-\gamma)>0,
\qquad
\lim_{t\to\infty}f(t)=1-\gamma>0.
\end{equation}
Moreover,
\begin{equation}
\label{eq:p0-f-derivative}
f'(t)=e^{-t}\left(2-t-\frac{A}{2}e^{t/2}\right),
\end{equation}
and the factor in parentheses is strictly decreasing.  Hence \(f\) has at most one critical point, and any such point is a maximum.  Its minimum on \([0,\infty)\) is attained at an endpoint, so \(f(t)>0\).  Thus \(K(t)>0\) for \(t>0\).

Differentiating under \eqref{eq:p0-laplace-kernel-representation} yields
\begin{equation}
\label{eq:p0-complete-monotonicity}
(-1)^r\frac{d^r}{dx^r}\bigl(-P_0''(x)\bigr)
=\int_0^\infty t^r e^{-xt}K(t)\,dt>0
\end{equation}
for all \(r\ge0\).  This proves strict complete monotonicity of \(-P_0''\), and the case \(r=0\) gives strict concavity of \(P_0\).
\end{proof}

Qi, Lim, and Nantomah also ask whether the higher-order products \(P_n(x)=\psi^{(n)}(x)\psi^{(n)}(1/x)\), \(n\in\N\), are convex, and more strongly whether \(P_n''\) is completely monotonic.  The following result refutes the stronger assertion.  It does not settle the weaker convexity question.

\begin{theorem}[APP-0008: Counterexample to complete monotonicity of higher-order polygamma product curvature]
\label{thm:pn-complete-monotonicity-counterexample}
Let
\begin{equation}
\label{eq:pn-product}
P_n(x)=\psi^{(n)}(x)\psi^{(n)}(1/x),\qquad n\in\N,\ x>0.
\end{equation}
Then \(P_n''\) is not completely monotonic for all \(n\in\N\).  More precisely, for every \(n\ge29\),
\begin{equation}
\label{eq:pn-high-order-third-positive}
P_n'''(2)>0,
\end{equation}
and, already at lower order,
\begin{equation}
\label{eq:p7-sixth-negative}
P_7^{(6)}(3)<0.
\end{equation}
\end{theorem}

\begin{proof}
For \(m\ge1\), define the sign-normalized polygamma function
\begin{equation}
\label{eq:sign-normalized-polygamma}
A_m(x)=(-1)^{m+1}\psi^{(m)}(x)
=m!\sum_{k=0}^{\infty}(x+k)^{-m-1}.
\end{equation}
Then \(A_m(x)>0\), \(A_m'(x)=-A_{m+1}(x)\), and \(P_n(x)=A_n(x)A_n(1/x)\).  Differentiating twice gives
\begin{equation}
\label{eq:pn-second-derivative-normalized}
\begin{aligned}
P_n''(x)=&
A_{n+2}(x)A_n(1/x)
-\frac{2}{x^2}A_{n+1}(x)A_{n+1}(1/x)\\
&-\frac{2}{x^3}A_n(x)A_{n+1}(1/x)
+\frac{1}{x^4}A_n(x)A_{n+2}(1/x).
\end{aligned}
\end{equation}
First prove \eqref{eq:pn-high-order-third-positive}.  Put \(p=n+1\), and normalize
\begin{equation}
\label{eq:qp-normalized-product}
Q_p(x)=\frac{P_n(x)}{(n!)^2}
=\sum_{m,\ell\ge0}
\left(\frac{x}{(x+m)(1+\ell x)}\right)^p.
\end{equation}
At \(x=2\), the term \((m,\ell)=(0,0)\) is constant, and the unique largest nonconstant base is \((m,\ell)=(1,0)\), equal to \(2/3\).  Write
\begin{equation}
\label{eq:qp-dominant-split}
Q_p(x)=1+\left(\frac{x}{x+1}\right)^p+R_p(x).
\end{equation}
For \(h(x)=x/(x+1)\), direct differentiation gives
\begin{equation}
\label{eq:h-third-derivative}
\frac{d^3}{dx^3}h(x)^p
=p\left(\frac{x}{x+1}\right)^p
\frac{p^2-6px-3p+6x^2+6x+2}{x^3(x+1)^3}.
\end{equation}
Therefore
\begin{equation}
\label{eq:h-third-at-two}
\frac{d^3}{dx^3}h(x)^p\bigg|_{x=2}
=\frac{p(p^2-15p+38)}{216}\left(\frac23\right)^p.
\end{equation}
For \(p\ge30\), this is at least
\begin{equation}
\label{eq:h-third-lower}
\frac{p^3}{432}\left(\frac23\right)^p.
\end{equation}

For a general summand
\begin{equation}
\label{eq:fmell-general}
F_{m,\ell}(x)=
\left(\frac{x}{(x+m)(1+\ell x)}\right)^p,
\end{equation}
let \(u_{m,\ell}=d(\log F_{m,\ell}^{1/p})/dx\).  At \(x=2\),
\begin{equation}
\label{eq:umell-bounds}
|u_{m,\ell}|\le\frac12,\qquad
|u_{m,\ell}'|\le\frac14,\qquad
|u_{m,\ell}''|\le\frac14.
\end{equation}
Since
\begin{equation}
\label{eq:fmell-third-log}
F_{m,\ell}'''=F_{m,\ell}\left(p^3u_{m,\ell}^3+3p^2u_{m,\ell}u_{m,\ell}'+pu_{m,\ell}''\right),
\end{equation}
we have \(|F_{m,\ell}'''(2)|\le p^3F_{m,\ell}(2)\) for \(p\ge2\).  The remaining mass in \(R_p\) is bounded by
\begin{equation}
\label{eq:rp-mass-bound}
\sum_{(m,\ell)\ne(0,0),(1,0)}F_{m,\ell}(2)<8\cdot2^{-p}
\qquad(p\ge3).
\end{equation}
Indeed, the \(m\ge2,\ell=0\) part is at most \(3\cdot2^{-p}\), while the \(\ell\ge1\) part is at most \(2\cdot2^{-p}(1+(2/3)^p+3\cdot2^{-p})\).
Combining \eqref{eq:h-third-lower} and \eqref{eq:rp-mass-bound},
\begin{equation}
\label{eq:qp-third-lower}
Q_p'''(2)\ge
p^3 2^{-p}
\left[
\frac1{432}\left(\frac43\right)^p-8
\right].
\end{equation}
Because \((4/3)^{30}>432\cdot8\), \eqref{eq:qp-third-lower} is positive for \(p\ge30\).  Hence \(P_n'''(2)=(n!)^2Q_{n+1}'''(2)>0\) for \(n\ge29\).  Complete monotonicity of \(P_n''\) would require \((P_n'')'=P_n'''\le0\), so \(P_n''\) is not completely monotonic for \(n\ge29\).

We now record the lower-order obstruction \eqref{eq:p7-sixth-negative}.  For bookkeeping, put
\begin{equation}
\label{eq:tpij-definition}
T_{p,i,j}(x)=x^{-p}A_i(x)A_j(1/x).
\end{equation}
The identity \(A_m'=-A_{m+1}\) gives the exact recurrence
\begin{equation}
\label{eq:tpij-derivative}
\frac{d}{dx}T_{p,i,j}(x)
=-pT_{p+1,i,j}(x)-T_{p,i+1,j}(x)+T_{p+2,i,j+1}(x).
\end{equation}
Starting with \eqref{eq:pn-second-derivative-normalized}, applying \eqref{eq:tpij-derivative} four times with \(n=7\) expresses \(P_7^{(6)}(3)\) as an exact finite linear combination of \(22\) terms \(T_{p,i,j}(3)\).

The sign of this finite expression can be certified by rational intervals.  For rational \(x>0\) and \(N\ge1\), monotone integral comparison in \eqref{eq:sign-normalized-polygamma} gives
\begin{equation}
\label{eq:am-rational-lower}
m!\left(\sum_{k=0}^{N-1}(x+k)^{-m-1}+\frac{(x+N)^{-m}}{m}\right)
\le A_m(x),
\end{equation}
and
\begin{equation}
\label{eq:am-rational-upper}
A_m(x)\le
m!\left(\sum_{k=0}^{N-1}(x+k)^{-m-1}+\frac{(x+N)^{-m}}{m}+(x+N)^{-m-1}\right).
\end{equation}
Using \(N=80\) in \eqref{eq:am-rational-lower} and \eqref{eq:am-rational-upper} for every \(A_m(3)\) and \(A_m(1/3)\) appearing in the \(22\)-term expression, and choosing upper or lower endpoints according to the sign of each coefficient, yields
\begin{equation}
\label{eq:p7-sixth-certified-lower}
-74998.154649422399828421756528409981771664439555587
<P_7^{(6)}(3),
\end{equation}
and
\begin{equation}
\label{eq:p7-sixth-certified-upper}
P_7^{(6)}(3)
<-74997.331312439140215056900814096498575289999496021.
\end{equation}
Thus \(P_7^{(6)}(3)<0\).  If \(P_7''\) were completely monotonic, then
\begin{equation}
\label{eq:p7-cm-required-sign}
(-1)^4\frac{d^4}{dx^4}P_7''(3)=P_7^{(6)}(3)\ge0,
\end{equation}
contradicting \eqref{eq:p7-sixth-certified-upper}.  Hence the proposed complete-monotonicity strengthening for all \(n\in\N\) is false in both an infinite high-order family and already at \(n=7\).
\end{proof}

\begin{proposition}[Zeta-law calculus]
\label{prop:zeta-law-calculus}
Let \(\beta>1\), and let \(N\) have law \(\rho_\beta\).  For every \(r\ge0\),
\begin{equation}
\label{eq:zeta-law-moment}
\E_\beta[N^{-r}]=\frac{\zeta(\beta+r)}{\zeta(\beta)}.
\end{equation}
Moreover, for \(\alpha,\beta>1\),
\begin{equation}
\label{eq:zeta-law-relative-entropy}
D(\rho_\alpha\Vert\rho_\beta)
=(\beta-\alpha)\E_\alpha[\log N]
+\log\zeta(\beta)-\log\zeta(\alpha).
\end{equation}
\end{proposition}

\begin{proof}
The moment identity follows by direct summation:
\begin{equation}
\label{eq:zeta-law-moment-proof}
\E_\beta[N^{-r}]
=\sum_{n=1}^{\infty}n^{-r}\frac{n^{-\beta}}{\zeta(\beta)}
=\frac{\zeta(\beta+r)}{\zeta(\beta)}.
\end{equation}
For relative entropy, substitute \eqref{eq:zeta-law}:
\begin{equation}
\label{eq:zeta-law-relative-entropy-proof}
D(\rho_\alpha\Vert\rho_\beta)
=\sum_{n=1}^{\infty}\rho_\alpha(n)
\left((\beta-\alpha)\log n+\log\zeta(\beta)-\log\zeta(\alpha)\right),
\end{equation}
which is \eqref{eq:zeta-law-relative-entropy}.
\end{proof}

\section{Tail zeta laws and inverse-tail floors}

\begin{definition}[Tail zeta partition function]
\label{def:tail-zeta-partition}
For \(s>1\) and \(n\in\N\), define the tail zeta partition function by
\begin{equation}
\label{eq:tail-zeta-partition}
Z_{s,n}=\zeta_n(s)=\sum_{k=n}^{\infty}k^{-s}.
\end{equation}
The corresponding tail zeta law is the probability distribution on \(\{n,n+1,\ldots\}\) given by
\begin{equation}
\label{eq:tail-zeta-law}
\rho_{s,n}(k)=\frac{k^{-s}}{\zeta_n(s)},\qquad k\ge n.
\end{equation}
\end{definition}

Exact formulas for integer parts of reciprocal zeta tails have been studied for small integer exponents.  Known formulas include \(s=2,3,4,5\), with a sufficiently-large-\(n\) formula for \(s=6\); Kim and Song record that no analogous formula was known for integer \(s>6\) at the time of their work \cite{kim-song-tail}.  The 2024 asymptotic inverse-tail results of Pan and Wu provide related context \cite{pan-wu-tail}.  The next two theorems give exact formulas at \(s=7\) and \(s=8\).  The first uses a single rational telescoping enclosure, while the second needs adjacent corrected asymptotic truncations.

\begin{theorem}[APP-0005: Exact inverse-tail floor formula at s=7]
\label{thm:tail-s7-floor}
Let
\begin{equation}
\label{eq:tail-s7}
T_7(n)=\zeta_n(7)=\sum_{k=n}^{\infty}\frac1{k^7}.
\end{equation}
Define
\begin{equation}
\label{eq:tail-s7-q}
Q(n)=120n^6-360n^5+660n^4-720n^3+354n^2-54n+375,
\qquad
P(n)=\frac{Q(n)}{20}.
\end{equation}
Then, for every \(n\ge28\),
\begin{equation}
\label{eq:tail-s7-floor}
\left\lfloor T_7(n)^{-1}\right\rfloor=\left\lfloor P(n)\right\rfloor.
\end{equation}
For \(1\le n\le27\), the exact values are
\begin{equation}
\label{eq:tail-s7-finite-table}
\begin{array}{c|r@{\qquad}c|r@{\qquad}c|r}
n&\lfloor T_7(n)^{-1}\rfloor&n&\lfloor T_7(n)^{-1}\rfloor&n&\lfloor T_7(n)^{-1}\rfloor\\
\hline
1&0&10&4495760&19&241765529\\
2&119&11&8167814&20&331399044\\
3&1862&12&14061542&21&447175152\\
4&12573&13&23143975&22&594869693\\
5&54069&14&36668777&23&781167220\\
6&175597&15&56228085&24&1013751716\\
7&471118&16&83808666&25&1301401638\\
8&1100904&17&121852400&26&1654089273\\
9&2317459&18&173321080&27&2083084422
\end{array}.
\end{equation}
\end{theorem}

\begin{proof}
The polynomial \(Q\) is positive on \(\N\), because
\begin{equation}
\label{eq:tail-s7-q-positive}
\begin{aligned}
Q(n)
&=120(n-1)^6+360(n-1)^5+660(n-1)^4+720(n-1)^3\\
&\qquad +354(n-1)^2+54(n-1)+375.
\end{aligned}
\end{equation}
For \(k\ge1\), direct expansion gives
\begin{equation}
\label{eq:tail-s7-lower-difference}
\frac1{k^7}
-\left(\frac{20}{Q(k)}-\frac{20}{Q(k+1)}\right)
=
\frac{9(60284k^4+29176k^2+15625)}
{k^7Q(k)Q(k+1)}>0.
\end{equation}
Summing \eqref{eq:tail-s7-lower-difference} over \(k\ge n\) yields
\begin{equation}
\label{eq:tail-s7-lower-sum}
T_7(n)>\frac{20}{Q(n)}=\frac1{P(n)},
\end{equation}
and therefore \(T_7(n)^{-1}<P(n)\).

For \(k\ge28\), another expansion gives
\begin{equation}
\label{eq:tail-s7-upper-difference}
\left(\frac{20}{Q(k)-3}-\frac{20}{Q(k+1)-3}\right)
-\frac1{k^7}
=
\frac{36(20k^6-14961k^4-7235k^2-3844)}
{k^7(Q(k)-3)(Q(k+1)-3)}.
\end{equation}
The numerator in \eqref{eq:tail-s7-upper-difference} is positive for \(k\ge28\).  Indeed, writing \(k=m+28\), it becomes
\begin{equation}
\label{eq:tail-s7-upper-positive}
20m^6+3360m^5+220239m^4+7105168m^3
+114013021m^2+751143512m+436261580,
\end{equation}
which has positive coefficients for \(m\ge0\).  Hence, for \(n\ge28\),
\begin{equation}
\label{eq:tail-s7-upper-sum}
T_7(n)<\frac{20}{Q(n)-3}=\frac1{P(n)-3/20}.
\end{equation}
Combining \eqref{eq:tail-s7-lower-sum} and \eqref{eq:tail-s7-upper-sum} gives
\begin{equation}
\label{eq:tail-s7-inverse-window}
P(n)-\frac3{20}<T_7(n)^{-1}<P(n).
\end{equation}

Modulo \(20\),
\begin{equation}
\label{eq:tail-s7-q-mod}
Q(n)\equiv 14n^2-14n+15\pmod {20}.
\end{equation}
Since \(n(n-1)\pmod {10}\in\{0,2,6\}\), this implies
\begin{equation}
\label{eq:tail-s7-q-residues}
Q(n)\pmod {20}\in\{3,15,19\}.
\end{equation}
Thus the fractional part of \(P(n)\) is one of \(3/20,15/20,19/20\).  If \(a=\lfloor P(n)\rfloor\), then \(P(n)-3/20\ge a\), and \eqref{eq:tail-s7-inverse-window} gives
\begin{equation}
\label{eq:tail-s7-floor-trap}
a<T_7(n)^{-1}<a+1.
\end{equation}
This proves \eqref{eq:tail-s7-floor} for \(n\ge28\).

For the finite range, \(T_7(1)>1\), so \(0<T_7(1)^{-1}<1\).  For \(2\le n\le27\), let \(M_n\) denote the corresponding value in \eqref{eq:tail-s7-finite-table}.  With \(N=1500\), exact rational arithmetic verifies
\begin{equation}
\label{eq:tail-s7-finite-lower}
\sum_{k=n}^{N}\frac1{k^7}>\frac1{M_n+1}
\end{equation}
and
\begin{equation}
\label{eq:tail-s7-finite-upper}
\sum_{k=n}^{N}\frac1{k^7}+\frac1{6N^6}<\frac1{M_n}.
\end{equation}
The monotone integral estimate
\begin{equation}
\label{eq:tail-s7-integral-tail}
\sum_{k=N+1}^{\infty}\frac1{k^7}<\int_N^\infty x^{-7}\,dx=\frac1{6N^6}
\end{equation}
then gives
\begin{equation}
\label{eq:tail-s7-finite-window}
\frac1{M_n+1}<T_7(n)<\frac1{M_n}.
\end{equation}
Therefore \(M_n<T_7(n)^{-1}<M_n+1\), proving \eqref{eq:tail-s7-finite-table}.  The finite rational certificate is stored with the source as a reproducible calculation.
\end{proof}

\begin{theorem}[APP-0006: Exact inverse-tail floor formula at s=8]
\label{thm:tail-s8-floor}
Let
\begin{equation}
\label{eq:tail-s8}
T_8(n)=\zeta_n(8)=\sum_{k=n}^{\infty}\frac1{k^8}.
\end{equation}
Define
\begin{align}
\label{eq:tail-s8-a}
A_8(n)
={}&7n^7-\frac{49}{2}n^6+\frac{637}{12}n^5-\frac{1715}{24}n^4
+\frac{7399}{144}n^3-\frac{1715}{96}n^2+\frac{69685}{1728}n-\frac{65611}{3456}\notag\\
&-\frac{23764069}{103680n}
-\frac{23764069}{207360n^2}
+\frac{2102131409}{1244160n^3}
+\frac{2149659547}{829440n^4}\notag\\
&-\frac{203290800013}{14929920n^5}
-\frac{1146003910541}{29859840n^6}
+\frac{21587525169497}{179159040n^7}.
\end{align}
Then, for every \(n\ge6\),
\begin{equation}
\label{eq:tail-s8-floor}
\left\lfloor T_8(n)^{-1}\right\rfloor=\left\lfloor A_8(n)\right\rfloor.
\end{equation}
For \(1\le n\le5\), the exact values are
\begin{equation}
\label{eq:tail-s8-finite-table}
\begin{array}{c|rrrrr}
n&1&2&3&4&5\\
\hline
\lfloor T_8(n)^{-1}\rfloor&0&245&5844&53503&291407
\end{array}.
\end{equation}
\end{theorem}

\begin{proof}
Set
\begin{equation}
\label{eq:tail-s8-b}
B_8(n)=A_8(n)+\frac{200335112892761}{358318080n^8}.
\end{equation}
For \(C=A_8\) or \(B_8\), write \(R_C(k)=1/C(k)\).  Exact denominator clearing gives the two telescoping inequalities
\begin{equation}
\label{eq:tail-s8-lower-inverse-difference}
R_{A_8}(k)-R_{A_8}(k+1)-\frac1{k^8}>0\qquad(k\ge3)
\end{equation}
and
\begin{equation}
\label{eq:tail-s8-upper-inverse-difference}
\frac1{k^8}-\left(R_{B_8}(k)-R_{B_8}(k+1)\right)>0\qquad(k\ge1).
\end{equation}
In \eqref{eq:tail-s8-lower-inverse-difference}, the cleared numerator is a polynomial in \(k-3\) with all coefficients positive; in \eqref{eq:tail-s8-upper-inverse-difference}, the cleared numerator is a polynomial in \(k-1\) with all coefficients positive.  The exact coefficient verification is recorded in the rational certificate accompanying the manuscript.

Summing \eqref{eq:tail-s8-lower-inverse-difference} over \(k\ge n\) gives
\begin{equation}
\label{eq:tail-s8-lower-inverse-sum}
T_8(n)<R_{A_8}(n)=\frac1{A_8(n)}
\end{equation}
for \(n\ge3\), and hence \(A_8(n)<T_8(n)^{-1}\).  Summing \eqref{eq:tail-s8-upper-inverse-difference} gives
\begin{equation}
\label{eq:tail-s8-upper-inverse-sum}
T_8(n)>R_{B_8}(n)=\frac1{B_8(n)},
\end{equation}
and hence \(T_8(n)^{-1}<B_8(n)\).  Therefore, for \(n\ge6\),
\begin{equation}
\label{eq:tail-s8-inverse-window}
A_8(n)<T_8(n)^{-1}<B_8(n).
\end{equation}

It remains to see that no integer lies between \(A_8(n)\) and \(B_8(n)\).  Exact integer arithmetic verifies
\begin{equation}
\label{eq:tail-s8-floor-gap-finite}
\lfloor A_8(n)\rfloor=\lfloor B_8(n)\rfloor\qquad(6\le n\le10^6).
\end{equation}
For \(n>10^6\), let \(P_8(n)\) be the polynomial part of \(A_8(n)\), through the constant term.  The numerator of \(P_8(n)\) modulo \(3456\) runs through the odd residue classes, so \(P_8(n)\) is at distance at least \(1/3456\) from the next integer.  The correction \(A_8(n)-P_8(n)\) is negative and has absolute value \(<1/3456\), while
\begin{equation}
\label{eq:tail-s8-gap-small}
0<B_8(n)-A_8(n)<\frac1{3456}.
\end{equation}
Thus \(B_8(n)\) is still below the next integer after \(A_8(n)\).  Combining this with \eqref{eq:tail-s8-inverse-window} proves \eqref{eq:tail-s8-floor}.

For \(1\le n\le5\), \eqref{eq:tail-s8-finite-table} is certified by exact rational partial sums.  With \(N=300\), the displayed value \(M_n\) satisfies, for \(2\le n\le5\),
\begin{equation}
\label{eq:tail-s8-finite-lower}
\sum_{k=n}^{N}\frac1{k^8}>\frac1{M_n+1}
\end{equation}
and
\begin{equation}
\label{eq:tail-s8-finite-upper}
\sum_{k=n}^{N}\frac1{k^8}+\frac1{7N^7}<\frac1{M_n}.
\end{equation}
Since
\begin{equation}
\label{eq:tail-s8-integral-tail}
\sum_{k=N+1}^{\infty}\frac1{k^8}<\int_N^\infty x^{-8}\,dx=\frac1{7N^7},
\end{equation}
we get \(M_n<T_8(n)^{-1}<M_n+1\).  For \(n=1\), \(T_8(1)>1\), so the floor is \(0\).
\end{proof}

\section{Modular successor entropy}

\begin{definition}[Modular residue distribution and successor entropy]
\label{def:modular-residue-successor}
For \(q\ge2\), define the residue distribution of \(\rho_\beta\) modulo \(q\) by
\begin{equation}
\label{eq:modular-residue-distribution}
\mu_{q,\beta}(a)=
\sum_{\substack{n\ge1\\ n\equiv a\pmod q}}\rho_\beta(n),
\qquad a\in\Z/q\Z.
\end{equation}
Writing \(a+1\) cyclically in \(\Z/q\Z\), define
\begin{equation}
\label{eq:modular-successor-entropy}
B_q(\beta)=
\sum_{a\in\Z/q\Z}
\mu_{q,\beta}(a)\log\frac{\mu_{q,\beta}(a)}{\mu_{q,\beta}(a+1)}.
\end{equation}
\end{definition}

\begin{theorem}[Zeta-law successor entropy and modular resolution]
\label{thm:successor-entropy}
For every \(\beta>1\), define
\begin{equation}
\label{eq:successor-entropy}
\Sigma(\beta)=\sum_{n=1}^{\infty}\rho_\beta(n)
\log\frac{\rho_\beta(n)}{\rho_\beta(n+1)}.
\end{equation}
Then
\begin{equation}
\label{eq:successor-entropy-series}
\Sigma(\beta)=
\frac{\beta}{\zeta(\beta)}
\sum_{k=1}^{\infty}\frac{(-1)^{k+1}}{k}\zeta(\beta+k),
\end{equation}
and
\begin{equation}
\label{eq:successor-entropy-modular-resolution}
\Sigma(\beta)=\sup_{q\ge2}B_q(\beta)=\lim_{q\to\infty}B_q(\beta).
\end{equation}
\end{theorem}

\begin{proof}
Since
\begin{equation}
\label{eq:successor-ratio}
\frac{\rho_\beta(n)}{\rho_\beta(n+1)}
=\left(1+\frac1n\right)^\beta,
\end{equation}
we have
\begin{equation}
\label{eq:successor-entropy-log-sum}
\Sigma(\beta)=\frac{\beta}{\zeta(\beta)}
\sum_{n=1}^{\infty}n^{-\beta}\log\left(1+\frac1n\right).
\end{equation}
Using
\begin{equation}
\label{eq:log-series}
\log(1+x)=\sum_{k=1}^{\infty}\frac{(-1)^{k+1}}{k}x^k
\qquad(0<x\le1)
\end{equation}
with Abel limiting at \(x=1\), and summing against \(n^{-\beta}\), gives \eqref{eq:successor-entropy-series}.

For the modular identity, fix \(q\ge2\).  The log-sum inequality gives, for each residue class \(a\),
\begin{equation}
\label{eq:modular-log-sum}
\sum_{\substack{n\ge1\\ n\equiv a\pmod q}}
\rho_\beta(n)\log\frac{\rho_\beta(n)}{\rho_\beta(n+1)}
\ge
\mu_{q,\beta}(a)
\log\frac{\mu_{q,\beta}(a)}
{\sum_{n\equiv a\pmod q}\rho_\beta(n+1)}.
\end{equation}
For \(a\ne0\), the denominator is \(\mu_{q,\beta}(a+1)\).  For \(a=0\), it is \(\mu_{q,\beta}(1)-\rho_\beta(1)\le\mu_{q,\beta}(1)\).  Summing \eqref{eq:modular-log-sum} over \(a\) gives
\begin{equation}
\label{eq:successor-dominates-modular}
\Sigma(\beta)\ge B_q(\beta)
\qquad(q\ge2).
\end{equation}

Conversely, fix \(M\in\N\).  If \(q>M+1\), then for \(1\le a\le M+1\),
\begin{equation}
\label{eq:residue-isolation}
\mu_{q,\beta}(a)=\rho_\beta(a)+O_{\beta,M}(q^{-\beta}),
\end{equation}
and \(\mu_{q,\beta}(0)=q^{-\beta}\).  All cyclic terms in \(B_q(\beta)\) are nonnegative except possibly the boundary term, which satisfies
\begin{equation}
\label{eq:boundary-term}
\mu_{q,\beta}(0)\log\frac{\mu_{q,\beta}(0)}{\mu_{q,\beta}(1)}
=O_\beta(q^{-\beta}\log q).
\end{equation}
Therefore
\begin{equation}
\label{eq:modular-liminf}
\liminf_{q\to\infty}B_q(\beta)
\ge
\sum_{n=1}^{M}\rho_\beta(n)\log\frac{\rho_\beta(n)}{\rho_\beta(n+1)}.
\end{equation}
Letting \(M\to\infty\) gives \(\liminf_qB_q(\beta)\ge\Sigma(\beta)\), which together with \eqref{eq:successor-dominates-modular} proves \eqref{eq:successor-entropy-modular-resolution}.
\end{proof}

\begin{corollary}[Prime-modulus Dirichlet L-resolution]
\label{cor:prime-dirichlet-resolution}
Let \(p\) be prime, let \(\chi\) range over Dirichlet characters modulo \(p\), and define
\begin{equation}
\label{eq:prime-dirichlet-shadow}
S_a(p,\beta)=\sum_{\chi\bmod p}\overline{\chi(a)}L(\beta,\chi)
\qquad(1\le a\le p-1).
\end{equation}
Then
\begin{equation}
\label{eq:prime-residue-formula}
\mu_{p,\beta}(a)=\frac{S_a(p,\beta)}{(p-1)\zeta(\beta)}
\quad(1\le a\le p-1),
\qquad
\mu_{p,\beta}(0)=p^{-\beta}.
\end{equation}
Consequently the prime-modulus functionals obtained by substituting \eqref{eq:prime-residue-formula} into \(B_p(\beta)\) satisfy
\begin{equation}
\label{eq:prime-functional-limit}
\Sigma(\beta)=
\lim_{\substack{p\to\infty\\ p\ {\rm prime}}}
F_p\bigl(\zeta(\beta),\{L(\beta,\chi)\}_{\chi\bmod p}\bigr)
=
\sup_{p\ {\rm prime}}
F_p\bigl(\zeta(\beta),\{L(\beta,\chi)\}_{\chi\bmod p}\bigr).
\end{equation}
\end{corollary}

\begin{proof}
For \(a\not\equiv0\pmod p\), character orthogonality gives
\begin{equation}
\label{eq:character-orthogonality}
\frac1{p-1}\sum_{\chi\bmod p}\overline{\chi(a)}\chi(n)
=
\begin{cases}
1,& n\equiv a\pmod p,\\
0,& n\not\equiv a\pmod p.
\end{cases}
\end{equation}
Therefore
\begin{equation}
\label{eq:dirichlet-residue-sum}
\sum_{\substack{n\ge1\\ n\equiv a\pmod p}}n^{-\beta}
=\frac1{p-1}\sum_{\chi\bmod p}\overline{\chi(a)}L(\beta,\chi).
\end{equation}
Dividing by \(\zeta(\beta)\) gives the nonzero residue formula, while
\begin{equation}
\label{eq:zero-residue-prime}
\mu_{p,\beta}(0)=\frac1{\zeta(\beta)}\sum_{m=1}^{\infty}(mp)^{-\beta}
=p^{-\beta}.
\end{equation}
The functional identity follows by substituting these residue formulas into \(B_p(\beta)\) and applying Theorem \ref{thm:successor-entropy} along the prime sequence.
\end{proof}

\section{Euler-score identities}

\begin{theorem}[Euler-score decomposition]
\label{thm:euler-score-decomposition}
For every \(\beta>1\),
\begin{equation}
\label{eq:euler-score-decomposition}
-\bigl(\log\zeta\bigr)'(\beta)
=\E_\beta[\log N]
=\sum_{d=1}^{\infty}\Lambda(d)\mu_{d,\beta}(0)
=\sum_{d=1}^{\infty}\frac{\Lambda(d)}{d^\beta}.
\end{equation}
Equivalently,
\begin{equation}
\label{eq:log-zeta-von-mangoldt}
-\frac{\zeta'(\beta)}{\zeta(\beta)}
=\sum_{d=1}^{\infty}\frac{\Lambda(d)}{d^\beta}.
\end{equation}
\end{theorem}

\begin{proof}
The von Mangoldt identity is
\begin{equation}
\label{eq:mangoldt-identity}
\log n=\sum_{d\mid n}\Lambda(d).
\end{equation}
By nonnegative summation,
\begin{equation}
\label{eq:euler-score-proof}
\E_\beta[\log N]
=\sum_{n\ge1}\rho_\beta(n)\sum_{d\mid n}\Lambda(d)
=\sum_{d\ge1}\Lambda(d)\Prob_\beta(d\mid N).
\end{equation}
The divisibility probability is
\begin{equation}
\label{eq:divisibility-probability}
\Prob_\beta(d\mid N)
=\frac1{\zeta(\beta)}\sum_{m\ge1}(dm)^{-\beta}
=d^{-\beta}.
\end{equation}
Substitution into \eqref{eq:euler-score-proof} proves \eqref{eq:euler-score-decomposition}; \eqref{eq:log-zeta-von-mangoldt} follows from \eqref{eq:first-free-derivative-proof}.
\end{proof}

\begin{proposition}[Finite Euler-score identity]
\label{prop:finite-euler-score}
For every \(A\subseteq\{1,\ldots,N\}\),
\begin{equation}
\label{eq:finite-euler-score}
\sum_{n\in A}\log n
=\sum_{d\le N}\Lambda(d)\sum_{m\le N/d}\1_A(md).
\end{equation}
\end{proposition}

\begin{proof}
Reverse the order of summation:
\begin{equation}
\label{eq:finite-euler-score-proof}
\sum_{d\le N}\Lambda(d)\sum_{m\le N/d}\1_A(md)
=\sum_{n\in A}\sum_{d\mid n}\Lambda(d)
=\sum_{n\in A}\log n.
\end{equation}
\end{proof}

\section{Curvature and Mellin-Planck tools}

\begin{lemma}[Uniform positive-axis curvature bound]
\label{lem:uniform-positive-axis-curvature}
For every \(s\ge3\),
\begin{equation}
\label{eq:uniform-curvature-bound}
(\log\zeta)''(s)=\Var_s(\log N)<\frac13.
\end{equation}
\end{lemma}

\begin{proof}
By Proposition \ref{prop:free-energy-derivatives},
\begin{equation}
\label{eq:curvature-first-bound}
(\log\zeta)''(s)
=\frac{\zeta''(s)}{\zeta(s)}
-\left(\frac{\zeta'(s)}{\zeta(s)}\right)^2
\le \frac{\zeta''(s)}{\zeta(s)}
\le \zeta''(s).
\end{equation}
For \(s\ge3\),
\begin{equation}
\label{eq:zeta-second-derivative-bound}
\zeta''(s)=\sum_{m=2}^{\infty}\frac{(\log m)^2}{m^s}
\le\sum_{m=2}^{\infty}\frac{(\log m)^2}{m^3}.
\end{equation}
The function \(x\mapsto(\log x)^2x^{-3}\) is decreasing for \(x\ge2\), so
\begin{equation}
\label{eq:curvature-integral-comparison}
\sum_{m=2}^{\infty}\frac{(\log m)^2}{m^3}
\le
\frac{(\log2)^2}{8}
+\int_2^\infty\frac{(\log x)^2}{x^3}\,dx.
\end{equation}
Since
\begin{equation}
\label{eq:curvature-integral-value}
\int_2^\infty\frac{(\log x)^2}{x^3}\,dx
=\frac{(\log2)^2+\log2+\frac12}{8},
\end{equation}
we obtain
\begin{equation}
\label{eq:curvature-final-bound}
(\log\zeta)''(s)
\le
\frac{2(\log2)^2+\log2+\frac12}{8}
<\frac13.
\end{equation}
\end{proof}

\begin{definition}[Mellin-Planck partition function]
\label{def:mellin-planck-partition}
For \(s>1\), define
\begin{equation}
\label{eq:mellin-planck-partition}
M(s)=\Gamma(s)\zeta(s)
=\int_0^\infty\frac{t^{s-1}}{e^t-1}\,dt.
\end{equation}
The corresponding probability law is
\begin{equation}
\label{eq:mellin-planck-law}
\nu_s(dt)=\frac{t^{s-1}}{(e^t-1)M(s)}\,dt.
\end{equation}
\end{definition}

\begin{proposition}[Mellin-Planck log-convexity]
\label{prop:mellin-planck-log-convexity}
For \(s>1\),
\begin{equation}
\label{eq:mellin-planck-log-convexity}
\frac{d^2}{ds^2}\log M(s)=\Var_{\nu_s}(\log t)\ge0.
\end{equation}
\end{proposition}

\begin{proof}
Differentiate the logarithm of the integral in \eqref{eq:mellin-planck-partition} under the integral sign.  The first derivative is the \(\nu_s\)-expectation of \(\log t\), and the second derivative is its variance:
\begin{equation}
\label{eq:mellin-planck-variance-proof}
\frac{d^2}{ds^2}\log M(s)
=\E_{\nu_s}[(\log t)^2]-\E_{\nu_s}[\log t]^2.
\end{equation}
\end{proof}

Bulboaca and Zayed proved several monotonicity criteria for functions of the form
\(u(s)/(\Gamma(s+\rho)\Gamma(s))\) and asked for best possible positive
values of \(\rho\) in their examples \cite{bulboaca-zayed-gamma}.  The next
application gives the sharp endpoint in the base reciprocal Gamma-product
case and the corresponding endpoint criterion for monotone logarithmic
derivatives.

\begin{theorem}[APP-0009: Sharp reciprocal Gamma-product monotonicity threshold]
\label{thm:gamma-product-sharp-threshold}
Let \(\psi=\Gamma'/\Gamma\), let \(\gamma\) be Euler's constant, and let
\(\rho_*>0\) be the unique solution of
\begin{equation}
\label{eq:gamma-product-rho-star}
\psi(1+\rho_*)=\gamma.
\end{equation}
For \(\rho>-1\), define
\begin{equation}
\label{eq:gamma-product-reciprocal}
\varphi_\rho(s)=\frac{1}{\Gamma(s+\rho)\Gamma(s)},\qquad s\ge1.
\end{equation}
Then \(\varphi_\rho\) is strictly decreasing on \([1,\infty)\) if and only if
\begin{equation}
\label{eq:gamma-product-threshold}
\rho\ge\rho_*.
\end{equation}
Equivalently,
\begin{equation}
\label{eq:gamma-product-W}
W_\rho(s)=\Gamma(s+\rho)\Gamma(s)
\end{equation}
is strictly increasing on \([1,\infty)\) if and only if
\eqref{eq:gamma-product-threshold} holds.  Numerically,
\begin{equation}
\label{eq:gamma-product-rho-star-numeric}
\rho_*=1.2583969670859318174106234224981693941\ldots .
\end{equation}

More generally, let \(u:[1,\infty)\to(0,\infty)\) be differentiable and suppose
\begin{equation}
\label{eq:gamma-product-J}
J(s)=\frac{u'(s)}{u(s)}
\end{equation}
is nonincreasing on \([1,\infty)\).  Then
\begin{equation}
\label{eq:gamma-product-general-Phi}
\Phi_{\rho,u}(s)=\frac{u(s)}{\Gamma(s+\rho)\Gamma(s)}
\end{equation}
is strictly decreasing on \([1,\infty)\) if and only if
\begin{equation}
\label{eq:gamma-product-general-threshold}
\psi(1+\rho)\ge\gamma+J(1).
\end{equation}
\end{theorem}

\begin{proof}
The digamma function is strictly increasing on \((0,\infty)\), since
\begin{equation}
\label{eq:digamma-positive-series}
\psi'(x)=\sum_{k=0}^{\infty}\frac{1}{(x+k)^2}>0.
\end{equation}
Hence \eqref{eq:gamma-product-rho-star} has a unique positive solution.

For \(s\ge1\), the logarithmic derivative of
\eqref{eq:gamma-product-W} is
\begin{equation}
\label{eq:gamma-product-log-derivative}
\frac{d}{ds}\log W_\rho(s)=\psi(s+\rho)+\psi(s).
\end{equation}
The right-hand side is strictly increasing in \(s\).  Its minimum on
\([1,\infty)\) is therefore
\begin{equation}
\label{eq:gamma-product-left-endpoint}
\psi(1+\rho)+\psi(1)=\psi(1+\rho)-\gamma.
\end{equation}
If \(\rho\ge\rho_*\), then \eqref{eq:gamma-product-left-endpoint} is
nonnegative, and it is positive for every \(s>1\).  Thus, for
\(1\le a<b\),
\begin{equation}
\label{eq:gamma-product-integrated-positive}
\log W_\rho(b)-\log W_\rho(a)
=\int_a^b\bigl(\psi(s+\rho)+\psi(s)\bigr)\,ds>0.
\end{equation}
So \(W_\rho\) is strictly increasing and \(\varphi_\rho=1/W_\rho\) is
strictly decreasing.  Conversely, if \(\rho<\rho_*\), then
\eqref{eq:gamma-product-left-endpoint} is negative, so by continuity
\(W_\rho\) decreases initially and \(\varphi_\rho\) increases initially.
This proves the base equivalence.

For the general statement,
\begin{equation}
\label{eq:gamma-product-general-log-derivative}
\frac{d}{ds}\log\Phi_{\rho,u}(s)
=J(s)-\psi(s+\rho)-\psi(s).
\end{equation}
The term \(J\) is nonincreasing by hypothesis, while
\(\psi(s+\rho)+\psi(s)\) is strictly increasing.  Therefore the logarithmic
derivative in \eqref{eq:gamma-product-general-log-derivative} is strictly
decreasing, and its maximum on \([1,\infty)\) is
\begin{equation}
\label{eq:gamma-product-general-left-endpoint}
J(1)-\psi(1+\rho)-\psi(1)
=J(1)-\psi(1+\rho)+\gamma.
\end{equation}
The function \(\Phi_{\rho,u}\) is strictly decreasing exactly when this maximum
is nonpositive, which is \eqref{eq:gamma-product-general-threshold}.  If the
inequality fails, the logarithmic derivative is positive near \(s=1\), so
monotonicity fails.
\end{proof}

Qi, Lim, and Nantomah asked whether the second derivative of \(\log\Gamma(x)+\log\Gamma(1/x)\) is completely monotonic on \((0,\infty)\) \cite{qi-lim-nantomah-polygamma}.  The next theorem answers this in the affirmative by a reciprocal-Weierstrass Laplace kernel.

\begin{theorem}[APP-0004: Complete monotonicity of reciprocal-Gamma curvature]
\label{thm:reciprocal-gamma-complete-monotonicity}
Let
\begin{equation}
\label{eq:reciprocal-gamma-h}
H(x)=\log\Gamma(x)+\log\Gamma(1/x),\qquad x>0.
\end{equation}
Then \(H''\) is strictly completely monotonic on \((0,\infty)\); that is,
\begin{equation}
\label{eq:reciprocal-gamma-complete-monotonicity}
(-1)^m\frac{d^m}{dx^m}H''(x)>0
\qquad(m\ge0,\ x>0).
\end{equation}
\end{theorem}

\begin{proof}
The Weierstrass product gives, for \(x>0\),
\begin{equation}
\label{eq:gamma-weierstrass-log}
\log\Gamma(x)
=-\log x-\gamma x-\sum_{n=1}^{\infty}
\left(\log\left(1+\frac{x}{n}\right)-\frac{x}{n}\right).
\end{equation}
After substituting \(1/x\) in \eqref{eq:gamma-weierstrass-log}, adding the two formulas, and differentiating twice on compact subintervals of \((0,\infty)\), the logarithmic second-derivative terms cancel and
\begin{equation}
\label{eq:reciprocal-gamma-second-series}
\begin{aligned}
H''(x)
&=\sum_{n=1}^{\infty}\frac1{(x+n)^2}
-\frac{2\gamma}{x^3}\\
&\qquad+
\sum_{n=1}^{\infty}
\left[
\frac1{(x+1/n)^2}
-\frac1{x^2}
+\frac{2}{nx^3}
\right].
\end{aligned}
\end{equation}
The elementary Laplace identities
\begin{equation}
\label{eq:reciprocal-gamma-laplace-identities}
\frac1{(x+a)^2}=\int_0^\infty t e^{-(x+a)t}\,dt,\qquad
\frac1{x^2}=\int_0^\infty t e^{-xt}\,dt,\qquad
\frac{2a}{x^3}=a\int_0^\infty t^2e^{-xt}\,dt
\end{equation}
show that, with \(a=1/n\),
\begin{equation}
\label{eq:reciprocal-gamma-defect-laplace}
\frac1{(x+1/n)^2}-\frac1{x^2}+\frac{2}{nx^3}
=\int_0^\infty e^{-xt}t\left(e^{-t/n}-1+\frac{t}{n}\right)\,dt.
\end{equation}
Also,
\begin{equation}
\label{eq:reciprocal-gamma-trigamma-laplace}
\sum_{n=1}^{\infty}\frac1{(x+n)^2}
=\int_0^\infty e^{-xt}\frac{t}{e^t-1}\,dt,
\end{equation}
and
\begin{equation}
\label{eq:reciprocal-gamma-euler-laplace}
-\frac{2\gamma}{x^3}
=\int_0^\infty e^{-xt}(-\gamma t^2)\,dt.
\end{equation}
Combining \eqref{eq:reciprocal-gamma-second-series}--\eqref{eq:reciprocal-gamma-euler-laplace} gives
\begin{equation}
\label{eq:reciprocal-gamma-kernel-representation}
H''(x)=\int_0^\infty e^{-xt}K(t)\,dt,
\end{equation}
where
\begin{equation}
\label{eq:reciprocal-gamma-kernel}
K(t)=\frac{t}{e^t-1}
+t\sum_{n=1}^{\infty}\left(e^{-t/n}-1+\frac{t}{n}\right)
-\gamma t^2.
\end{equation}

It remains to prove \(K(t)>0\).  Put
\begin{equation}
\label{eq:reciprocal-gamma-phi}
\phi(u)=e^{-u}-1+u.
\end{equation}
Then \(\phi\ge0\), \(\phi'=1-e^{-u}\ge0\), and \(u\mapsto\phi(t/u)\) is positive and decreasing on \([1,\infty)\).  Therefore
\begin{equation}
\label{eq:reciprocal-gamma-sum-integral}
\sum_{n=1}^{\infty}\phi(t/n)\ge\int_1^\infty \phi(t/u)\,du.
\end{equation}
With \(v=t/u\), integration by parts gives
\begin{equation}
\label{eq:reciprocal-gamma-integral-transform}
\int_1^\infty\phi(t/u)\,du
=t\int_0^t\frac{\phi(v)}{v^2}\,dv
=-\phi(t)+t\int_0^t\frac{1-e^{-v}}{v}\,dv.
\end{equation}
Using
\begin{equation}
\label{eq:reciprocal-gamma-e1-identity}
\int_0^t\frac{1-e^{-v}}{v}\,dv
=\gamma+\log t+E_1(t),
\qquad
E_1(t)=\int_t^\infty\frac{e^{-v}}{v}\,dv,
\end{equation}
we obtain
\begin{equation}
\label{eq:reciprocal-gamma-kernel-lower-bound}
\frac{K(t)}{t}
\ge
\frac1{e^t-1}-\phi(t)+t\log t+tE_1(t).
\end{equation}
Since \(E_1(t)>0\) and \(\log t\ge1-1/t\) for \(t>0\),
\begin{equation}
\label{eq:reciprocal-gamma-kernel-positive}
\frac{K(t)}{t}
>
\frac1{e^t-1}-\phi(t)+t-1
=\frac1{e^t-1}-e^{-t}
=\frac{e^{-t}}{e^t-1}>0.
\end{equation}
Thus \(K(t)>0\) for \(t>0\).  Differentiating under the Laplace integral in \eqref{eq:reciprocal-gamma-kernel-representation} gives
\begin{equation}
\label{eq:reciprocal-gamma-alternating-integral}
(-1)^m\frac{d^m}{dx^m}H''(x)
=\int_0^\infty t^m e^{-xt}K(t)\,dt>0.
\end{equation}
This proves \eqref{eq:reciprocal-gamma-complete-monotonicity}.
\end{proof}

\section{External zeta inequality proofs}

\begin{theorem}[APP-0002: Affirmative solution of Nantomah zeta positivity problem]
\label{thm:nantomah-positivity}
For every \(n\in\N\),
\begin{equation}
\label{eq:nantomah-positivity}
(n+2)\zeta(n+1)\zeta(n+3)
-(n+1)\zeta(n+2)^2
-\zeta(n+1)\zeta(n+2)>0.
\end{equation}
\end{theorem}

\begin{proof}
Put \(s=n+1\ge2\) and \(Z_s=\zeta(s)\).  The expression in \eqref{eq:nantomah-positivity} becomes
\begin{equation}
\label{eq:nantomah-ks}
K_s=(s+1)Z_sZ_{s+2}-sZ_{s+1}^2-Z_sZ_{s+1}.
\end{equation}
Let \(X(m)=1/m\) under \(\rho_s\).  By \eqref{eq:zeta-law-moment},
\begin{equation}
\label{eq:nantomah-moments}
\E_s[X]=\frac{Z_{s+1}}{Z_s},
\qquad
\E_s[X^2]=\frac{Z_{s+2}}{Z_s}.
\end{equation}
Thus
\begin{equation}
\label{eq:nantomah-moment-form}
\frac{K_s}{Z_s^2}
=(s+1)\E_s[X^2]-s\E_s[X]^2-\E_s[X]
=(s+1)\Var_s(X)+\E_s[X]^2-\E_s[X].
\end{equation}

Write \(a=\zeta(s)-1\), \(b=\zeta(s+1)-1\), and \(c=\zeta(s+2)-1\).  Expanding gives
\begin{align}
\label{eq:nantomah-lq-split}
K_s&=L_s+Q_s,&
L_s&=sa+(s+1)c-(2s+1)b,&
Q_s&=(s+1)ac-sb^2-ab.
\end{align}
The linear part is termwise positive:
\begin{equation}
\label{eq:nantomah-linear-positive}
L_s=
\sum_{m=2}^{\infty}
m^{-s}\frac{(m-1)(s(m-1)-1)}{m^2}
\ge
2^{-s}\frac{s-1}{4}.
\end{equation}
Since \(b\le a/2\) and \(c\ge0\),
\begin{equation}
\label{eq:nantomah-quadratic-correction}
Q_s\ge -sb^2-ab\ge-\frac{s+2}{4}a^2.
\end{equation}
For \(s\ge4\),
\begin{equation}
\label{eq:nantomah-tail-bound}
a=\sum_{m=2}^{\infty}m^{-s}
\le2^{-s}+\int_2^\infty x^{-s}\,dx
=2^{-s}\left(1+\frac2{s-1}\right).
\end{equation}
Equations \eqref{eq:nantomah-linear-positive}, \eqref{eq:nantomah-quadratic-correction}, and \eqref{eq:nantomah-tail-bound} imply
\begin{equation}
\label{eq:nantomah-lower-bound}
K_s\ge
\frac14\left[
2^{-s}(s-1)
-(s+2)2^{-2s}\left(1+\frac2{s-1}\right)^2
\right].
\end{equation}
It is therefore enough to show
\begin{equation}
\label{eq:nantomah-sufficient-inequality}
2^s(s-1)>(s+2)\left(1+\frac2{s-1}\right)^2.
\end{equation}
Equivalently,
\begin{equation}
\label{eq:nantomah-f-function}
F(s)=\frac{2^s(s-1)^3}{(s+2)(s+1)^2}>1.
\end{equation}
Now \(F(4)=72/25>1\), and
\begin{equation}
\label{eq:nantomah-f-monotone}
\frac{d}{ds}\log F(s)
=\log2+\frac3{s-1}-\frac1{s+2}-\frac2{s+1}>0
\qquad(s\ge4).
\end{equation}
Thus \(K_s>0\) for \(s\ge4\).

For \(s=2\),
\begin{equation}
\label{eq:nantomah-k2}
K_2=3\zeta(2)\zeta(4)-2\zeta(3)^2-\zeta(2)\zeta(3).
\end{equation}
Using \(\zeta(3)<5/4\), \(\zeta(2)=\pi^2/6\), and \(\zeta(4)=\pi^4/90\),
\begin{equation}
\label{eq:nantomah-k2-bound}
K_2>
\frac{\pi^6}{180}-\frac{5\pi^2}{24}-\frac{25}{8}
>
\frac{10415360504209}{180000000000000}>0.
\end{equation}
For \(s=3\), the estimates \(\zeta(3)>251/216\), \(\zeta(5)>8051/7776\), and \(\zeta(4)<13/12\) yield
\begin{equation}
\label{eq:nantomah-k3-bound}
K_3>
\frac{251}{216}
\left(4\cdot\frac{8051}{7776}-\frac{13}{12}\right)
-3\left(\frac{13}{12}\right)^2
=\frac{13783}{419904}>0.
\end{equation}
Hence \(K_s>0\) for all integers \(s\ge2\), proving the theorem.  This answers Nantomah's question \cite{nantomah}.
\end{proof}

\begin{theorem}[APP-0001: Alzer-Kwong convexity and concavity pattern for reciprocal zeta]
\label{thm:alzer-kwong}
Let \(F(x)=1/\zeta(x)\).  For every integer \(n\ge1\),
\begin{equation}
\label{eq:alzer-kwong-convex}
F''(x)>0\quad\text{on }(-4n,-4n+2),
\end{equation}
and
\begin{equation}
\label{eq:alzer-kwong-concave}
F''(x)<0\quad\text{on }(-4n-2,-4n).
\end{equation}
\end{theorem}

\begin{proof}
Put \(x=-u\), where \(u>2\).  The functional equation gives
\begin{equation}
\label{eq:zeta-functional-equation-negative-axis}
\zeta(-u)
=-2^{-u}\pi^{-u-1}\sin\left(\frac{\pi u}{2}\right)\Gamma(u+1)\zeta(u+1),
\end{equation}
and hence
\begin{equation}
\label{eq:reciprocal-zeta-negative-axis}
\frac1{\zeta(-u)}
=-\frac{2^u\pi^{u+1}}
{\Gamma(u+1)\zeta(u+1)\sin(\pi u/2)}.
\end{equation}
Define the positive function
\begin{equation}
\label{eq:alzer-g-function}
G(u)=
\frac{2^u\pi^{u+1}}
{\Gamma(u+1)\zeta(u+1)|\sin(\pi u/2)|}.
\end{equation}
On \(u\in(4n-2,4n)\), \(\sin(\pi u/2)<0\), so \(F(-u)=G(u)\).  On \(u\in(4n,4n+2)\), \(\sin(\pi u/2)>0\), so \(F(-u)=-G(u)\).  Since second derivatives are unchanged by \(x=-u\), it remains to prove \(G''(u)>0\).

Let \(\ell(u)=\log G(u)\).  Then
\begin{equation}
\label{eq:alzer-log-g}
\ell(u)=u\log2+(u+1)\log\pi-\log\Gamma(u+1)-\log\zeta(u+1)
-\log\left|\sin\left(\frac{\pi u}{2}\right)\right|.
\end{equation}
Differentiating twice gives
\begin{equation}
\label{eq:alzer-log-g-second}
\ell''(u)
=-\psi'(u+1)-(\log\zeta)''(u+1)
+\frac{\pi^2}{4}\csc^2\left(\frac{\pi u}{2}\right).
\end{equation}
With \(s=u+1>3\),
\begin{equation}
\label{eq:trigamma-bound}
\psi'(s)=\sum_{k=0}^{\infty}\frac1{(s+k)^2}
<
\int_{s-1}^{\infty}\frac{dt}{t^2}
=\frac1{s-1}\le\frac12.
\end{equation}
Combining \eqref{eq:uniform-curvature-bound}, \eqref{eq:trigamma-bound}, and \(\csc^2(\pi u/2)\ge1\), we get
\begin{equation}
\label{eq:alzer-log-convexity}
\ell''(u)>
\frac{\pi^2}{4}-\frac12-\frac13
=\frac{\pi^2}{4}-\frac56>0.
\end{equation}
Since \(G>0\),
\begin{equation}
\label{eq:alzer-g-second-positive}
G''(u)=G(u)(\ell''(u)+\ell'(u)^2)>0.
\end{equation}
The two signs in \eqref{eq:alzer-kwong-convex} and \eqref{eq:alzer-kwong-concave} now follow from \(F(-u)=G(u)\) on \(u\in(4n-2,4n)\) and \(F(-u)=-G(u)\) on \(u\in(4n,4n+2)\).  This proves the Alzer--Kwong sign pattern \cite{alzer-kwong}.
\end{proof}

\begin{theorem}[APP-0003: Generalized Holder inequality for Gamma zeta]
\label{thm:sroysang-holder}
Let \(m\ge2\), \(r\ge1\), \(x_1,\ldots,x_m>1\), and \(p_1,\ldots,p_m>1\) satisfy
\begin{equation}
\label{eq:sroysang-holder-condition}
\sum_{i=1}^{m}\frac1{p_i}=\frac1r.
\end{equation}
Define
\begin{equation}
\label{eq:sroysang-t}
T=r\sum_{i=1}^{m}\frac{x_i}{p_i}.
\end{equation}
Then \(T>1\) and
\begin{equation}
\label{eq:sroysang-generalized-holder}
\zeta(T)\le
\frac{\prod_{i=1}^{m}(\Gamma(x_i)\zeta(x_i))^{r/p_i}}{\Gamma(T)}.
\end{equation}
\end{theorem}

\begin{proof}
Set
\begin{equation}
\label{eq:sroysang-lambda}
\lambda_i=\frac{r}{p_i}.
\end{equation}
Then \(\lambda_i>0\), \(\sum_i\lambda_i=1\), and \(T=\sum_i\lambda_i x_i>1\).  Define
\begin{equation}
\label{eq:sroysang-q}
q_i=\frac{p_i}{r}.
\end{equation}
Then \(\sum_i1/q_i=1\), and the assumptions imply \(q_i>1\) for all \(i\); otherwise one term \(1/p_i\) would already be at least \(1/r\), leaving no room for the remaining positive terms in \eqref{eq:sroysang-holder-condition}.

For \(t>0\), let
\begin{equation}
\label{eq:sroysang-fi}
f_i(t)=\left(\frac{t^{x_i-1}}{e^t-1}\right)^{1/p_i}.
\end{equation}
Then
\begin{equation}
\label{eq:sroysang-product}
\left(\prod_{i=1}^{m}f_i(t)\right)^r
=\frac{t^{r\sum_i(x_i-1)/p_i}}{e^t-1}
=\frac{t^{T-1}}{e^t-1}.
\end{equation}
Consequently,
\begin{equation}
\label{eq:sroysang-lr-norm}
\left\|\prod_{i=1}^{m}f_i\right\|_{L^r(0,\infty)}^r
=\int_0^\infty\frac{t^{T-1}}{e^t-1}\,dt
=\Gamma(T)\zeta(T).
\end{equation}
Generalized Holder gives
\begin{equation}
\label{eq:sroysang-holder-step}
\left\|\prod_{i=1}^{m}f_i\right\|_{L^r}
\le
\prod_{i=1}^{m}\|f_i\|_{L^{p_i}}.
\end{equation}
Moreover,
\begin{equation}
\label{eq:sroysang-pi-norm}
\|f_i\|_{L^{p_i}}^{p_i}
=\int_0^\infty\frac{t^{x_i-1}}{e^t-1}\,dt
=\Gamma(x_i)\zeta(x_i).
\end{equation}
Combining \eqref{eq:sroysang-lr-norm}, \eqref{eq:sroysang-holder-step}, and \eqref{eq:sroysang-pi-norm}, then raising to the power \(r\), gives
\begin{equation}
\label{eq:sroysang-predivide}
\Gamma(T)\zeta(T)
\le
\prod_{i=1}^{m}(\Gamma(x_i)\zeta(x_i))^{r/p_i}.
\end{equation}
Division by \(\Gamma(T)\) proves \eqref{eq:sroysang-generalized-holder}.  This solves the generalized Holder problem in Sroysang's notation, where his \(\xi(s)\) is \(\zeta(s)\) for \(s>1\) \cite{sroysang}.
\end{proof}

\begin{remark}[Free-energy interpretation of the Sroysang inequality]
The same result follows from log-convexity of \(M(s)=\Gamma(s)\zeta(s)\).  With \(\lambda_i=r/p_i\) and \(T=\sum_i\lambda_ix_i\), Jensen's inequality for \(\log M\) gives
\begin{equation}
\label{eq:sroysang-jensen}
M(T)\le\prod_{i=1}^{m}M(x_i)^{\lambda_i}
=\prod_{i=1}^{m}(\Gamma(x_i)\zeta(x_i))^{r/p_i},
\end{equation}
which is equivalent to \eqref{eq:sroysang-generalized-holder}.
\end{remark}

\section{Laplace moment-ratio application}
\label{sec:laplace-moment-ratio-application}

The complete-monotonicity proofs above repeatedly convert special-function curvature questions into positivity statements for Laplace kernels.  The next lemma records a reusable moment-ratio consequence of that viewpoint.

\begin{lemma}[Laplace moment-ratio bridge]
\label{lem:laplace-moment-ratio-bridge}
Let \(\mu\) be a positive Borel measure on \([0,\infty)\), and assume that the tilted moments
\begin{equation}
\label{eq:laplace-moment-ratio-aj}
a_j(x)=\int_0^\infty t^j e^{-xt}\,d\mu(t)
\end{equation}
are finite and positive for the displayed indices.  Then, for every \(n\ge0\),
\begin{equation}
\label{eq:laplace-moment-ratio-Bn}
B_n(x)=\frac{a_{n+1}(x)}{a_n(x)a_{n+2}(x)}
\end{equation}
is strictly increasing in \(x\).
\end{lemma}

\begin{proof}
Put \(r_j(x)=a_{j+1}(x)/a_j(x)\).  Cauchy--Schwarz gives
\begin{equation}
\label{eq:laplace-moment-log-convexity}
a_{j+1}(x)^2\le a_j(x)a_{j+2}(x),
\end{equation}
so \(r_j(x)\le r_{j+1}(x)\).  Since \(a_j'(x)=-a_{j+1}(x)\), differentiating \eqref{eq:laplace-moment-ratio-Bn} logarithmically gives
\begin{equation}
\label{eq:laplace-moment-ratio-derivative}
\frac{d}{dx}\log B_n(x)
=-\frac{a_{n+2}}{a_{n+1}}
+\frac{a_{n+1}}{a_n}
+\frac{a_{n+3}}{a_{n+2}}
=r_n(x)-r_{n+1}(x)+r_{n+2}(x).
\end{equation}
The right-hand side is at least \(r_n(x)>0\), because \(r_{n+2}(x)\ge r_{n+1}(x)\).  Hence \(B_n\) is strictly increasing.
\end{proof}

Yin and Zhang define the Nielsen \(k\)-beta function by
\begin{equation}
\label{eq:nielsen-k-beta-definition}
\beta_k(x)=\int_0^1\frac{t^{x-1}}{1+t^k}\,dt
=\int_0^\infty\frac{e^{-xt}}{1+e^{-kt}}\,dt,
\qquad k>0,\ x>0.
\end{equation}
They ask for the monotonicity of the derivative ratio involving \(x\beta_k(x)\) \cite{yin-zhang-k-beta}.  The bridge lemma gives the exact parity law.

\begin{theorem}[APP-0010: Nielsen \(k\)-beta derivative-ratio monotonicity]
\label{thm:nielsen-k-beta-derivative-ratio}
Let \(k>0\), let \(n\ge0\), and put \(f_k(x)=x\beta_k(x)\).  Then
\begin{equation}
\label{eq:nielsen-k-beta-ratio}
R_{k,n}(x)=
\frac{f_k^{(n+1)}(x)}
{f_k^{(n)}(x)f_k^{(n+2)}(x)}
\end{equation}
is strictly increasing on \((0,\infty)\) when \(n\) is odd, and strictly decreasing on \((0,\infty)\) when \(n\) is even.
\end{theorem}

\begin{proof}
From \eqref{eq:nielsen-k-beta-definition},
\begin{equation}
\label{eq:nielsen-k-beta-x-integral}
f_k(x)=x\int_0^\infty \frac{e^{-xt}}{1+e^{-kt}}\,dt.
\end{equation}
Integrating by parts gives
\begin{equation}
\label{eq:nielsen-k-beta-laplace-density}
f_k(x)
=\frac12+
\int_0^\infty e^{-xt}
\frac{k e^{-kt}}{(1+e^{-kt})^2}\,dt.
\end{equation}
Thus \(f_k\) is a Laplace transform of a positive measure consisting of an atom \(1/2\) at \(0\) and a positive density on \((0,\infty)\).  Define
\begin{equation}
\label{eq:nielsen-k-beta-aj}
a_j(x)=(-1)^j f_k^{(j)}(x).
\end{equation}
For the needed indices these are positive tilted moments of the measure in \eqref{eq:nielsen-k-beta-laplace-density}.  Lemma~\ref{lem:laplace-moment-ratio-bridge} shows that
\begin{equation}
\label{eq:nielsen-k-beta-Bn-increasing}
\frac{a_{n+1}(x)}{a_n(x)a_{n+2}(x)}
\end{equation}
is strictly increasing.  But
\begin{equation}
\label{eq:nielsen-k-beta-sign}
R_{k,n}(x)=(-1)^{n+1}
\frac{a_{n+1}(x)}{a_n(x)a_{n+2}(x)}.
\end{equation}
Multiplication by \((-1)^{n+1}\) gives strict increase for odd \(n\) and strict decrease for even \(n\).  This resolves the Nielsen \(k\)-beta derivative-ratio problem with the precise parity convention.
\end{proof}

\section{Endpoint and parameter applications}
\label{sec:endpoint-parameter-applications}

The next three applications add one certified endpoint result and two parameter-monotonicity results to the complete-monotonicity layer.  Only fully solved source questions are labelled.

\subsection{Exact \(n=2\) beta-window endpoint}

For \(s>1\), write
\begin{equation}
\label{eq:q2-zs-definition}
Z_s(a)=\sum_{m=0}^\infty (m+a)^{-s}.
\end{equation}
In Qi--Lim--Nantomah's notation \cite{qi-lim-nantomah-polygamma}, put
\begin{equation}
\label{eq:q2-cp-definition}
C_2(x)=\psi''(x)+x\psi^{(3)}(x),
\qquad
P_2(x)=\psi''(x)\psi''(1/x),
\end{equation}
and define
\begin{equation}
\label{eq:q2-beta-window}
\mathcal I_2=
\{\beta\in\mathbb R:\ x^\beta C_2(x)-P_2(x)<0\text{ for all }x>0\}.
\end{equation}
The remaining endpoint is controlled on \(0<x<1\) by
\begin{equation}
\label{eq:q2-ratio-q}
R(x)=\frac{P_2(x)}{C_2(x)}
=\frac{2Z_3(x)Z_3(1/x)}{3xZ_4(x)-Z_3(x)},
\qquad
Q_2(x)=\frac{\log R(x)}{\log x}.
\end{equation}

\begin{theorem}[APP-0011: exact \(n=2\) beta-window endpoint]
\label{thm:q2-exact-endpoint}
Let
\begin{equation}
\label{eq:q2-J-bracket}
J=\left[\frac{287345}{1000000},\frac{287346}{1000000}\right].
\end{equation}
There is a unique \(\xi\in J\) at which \(Q_2\) attains its global supremum on \((0,1)\).  Hence, with \(L_2=Q_2(\xi)\),
\begin{equation}
\label{eq:q2-I2-exact}
\mathcal I_2=(L_2,3].
\end{equation}
\end{theorem}

\begin{proof}
The identity \(\psi^{(m)}(x)=(-1)^{m+1}m!Z_{m+1}(x)\) gives \eqref{eq:q2-ratio-q}.  Define
\begin{equation}
\label{eq:q2-lambda}
\Lambda(x)=
-\frac{3Z_4(x)}{Z_3(x)}
+\frac{3Z_4(1/x)}{x^2Z_3(1/x)}
-\frac{6Z_4(x)-12xZ_5(x)}{3xZ_4(x)-Z_3(x)}
\end{equation}
and
\begin{equation}
\label{eq:q2-g-function}
G(x)=x\log x\,\Lambda(x)-\log R(x).
\end{equation}
Differentiating \eqref{eq:q2-ratio-q} gives
\begin{equation}
\label{eq:q2-derivative}
Q_2'(x)=\frac{G(x)}{x(\log x)^2}.
\end{equation}

All inequalities below are rational interval certificates using the Hurwitz-zeta integral-tail enclosure and the atanh logarithm enclosure.  First,
\begin{equation}
\label{eq:q2-endpoint-signs}
G\!\left(\frac{287345}{1000000}\right)>0,
\qquad
G\!\left(\frac{287346}{1000000}\right)<0.
\end{equation}
Second, on
\[
B=\left[\frac{1409}{5000},\frac{293}{1000}\right],
\]
the same enclosure proves
\begin{equation}
\label{eq:q2-h-positive}
H(x)=\Lambda(x)+x\Lambda'(x)>0.
\end{equation}
Since \(G'(x)=\log x\,H(x)\) and \(\log x<0\) on \(B\), \(G\) is strictly decreasing on \(B\).  Together with \eqref{eq:q2-endpoint-signs}, this gives a unique zero \(\xi\in J\), and \eqref{eq:q2-derivative} shows that \(Q_2\) is increasing before \(\xi\) and decreasing after \(\xi\) on \(B\).

It remains only to exclude points outside the compact bracket.  The endpoint certificate gives a rational lower witness \(q_J<Q_2(\xi)\) with
\begin{equation}
\label{eq:q2-qJ-lower}
\frac{115727}{50000}<q_J.
\end{equation}
On \((0,9/50]\) the scaled near-zero bound gives \(Q_2(x)<23/10<q_J\).  On \([9/10,1)\), a finite rational subdivision proves \(R(x)>1\), hence \(Q_2(x)<0<q_J\).  On
\[
\left[\frac9{50},\frac{1409}{5000}\right]
\cup
\left[\frac{293}{1000},\frac9{10}\right],
\]
a finite rational cover proves \(Q_2(x)<q_J\): the left bridge uses
\[
R(x)>Z_3(1/x)>x^{1157/500},
\]
and the remaining intervals use the direct comparison
\[
\log R(x)>\frac{115727}{50000}\log x.
\]
Thus no point outside \(B\) reaches the value at \(\xi\), so \(L_2=Q_2(\xi)\).

The source reduction gives the upper endpoint \(3\) and the convention that the lower endpoint is excluded when equality occurs.  At \(\beta=L_2\), equality occurs at \(x=\xi\); therefore the exact admissible set is \((L_2,3]\).
\end{proof}

\subsection{Baricz \(V_q\) parameter log-convexity}

Baricz records the question whether, for fixed \(x>0\), the function
\begin{equation}
\label{eq:baricz-vq-definition}
V_q(x)=\frac{2e^{x^2}}{\Gamma(q+1)}
\int_x^\infty e^{-t^2}(t^2-x^2)^q\,dt,
\qquad q>-1,
\end{equation}
is log-convex as a function of \(q\) \cite{baricz-turan-thesis}.

\begin{theorem}[APP-0012: Baricz \(V_q\) strict \(q\)-log-convexity]
\label{thm:baricz-vq-strict-logconvexity}
For each fixed \(x>0\), the map \(q\mapsto V_q(x)\) is strictly log-convex on \((-1,\infty)\).
\end{theorem}

\begin{proof}
Put \(u=t^2-x^2\).  Then
\begin{equation}
\label{eq:baricz-u-normal-form}
V_q(x)=\frac1{\Gamma(q+1)}
\int_0^\infty e^{-u}u^q(u+x^2)^{-1/2}\,du.
\end{equation}
Using
\[
r^{-1/2}=\frac1{\sqrt\pi}\int_0^\infty s^{-1/2}e^{-rs}\,ds
\]
and Tonelli's theorem, we obtain
\begin{equation}
\label{eq:baricz-laplace-normal-form}
V_q(x)=\frac1{\sqrt\pi}\int_0^\infty
s^{-1/2}e^{-x^2s}(1+s)^{-(q+1)}\,ds.
\end{equation}
Let \(a=q+1>0\).  Differentiation under the integral gives
\begin{equation}
\label{eq:baricz-complete-monotonicity}
(-1)^n\frac{d^n}{da^n}V_{a-1}(x)
=\frac1{\sqrt\pi}\int_0^\infty
(\log(1+s))^n s^{-1/2}e^{-x^2s}(1+s)^{-a}\,ds\ge0.
\end{equation}
Equivalently, after the change of variables \(y=\log(1+s)\), \(V_{a-1}(x)\) is a positive Laplace transform in \(a\).  The representing measure is not concentrated at a single point, because the density in \eqref{eq:baricz-laplace-normal-form} is positive on \(s>0\).  Strict Holder inequality for two different exponents therefore gives strict log-convexity in \(a\), and the shift \(a=q+1\) gives the stated strict log-convexity in \(q\).
\end{proof}

\subsection{Yin \((p,k)\)-digamma sharp necessity}

Yin proved sufficiency of \(\alpha\le1\) for complete monotonicity of
\begin{equation}
\label{eq:yin-pk-delta}
\delta_{p,k,\alpha}(x)
=x^\alpha\left[
\frac1k\log\frac{pkx}{x+k(p+1)}-\psi_{p,k}(x)
\right],
\qquad p\in\mathbb N,\ k>0,
\end{equation}
and asked whether this condition is also necessary \cite{yin-pk-digamma}.

\begin{theorem}[APP-0013: Yin \((p,k)\)-digamma sharp \(\alpha\)-necessity]
\label{thm:yin-pk-alpha-necessity}
If \(\delta_{p,k,\alpha}\) is completely monotone on \((0,\infty)\), then \(\alpha\le1\).  Combined with Yin's sufficiency theorem, the sharp condition is \(\alpha\le1\).
\end{theorem}

\begin{proof}
The source representation is
\begin{equation}
\label{eq:yin-pk-digamma-formula}
\psi_{p,k}(x)=\frac1k\log(pk)-\sum_{n=0}^p\frac1{x+nk}.
\end{equation}
Therefore the bracket in \eqref{eq:yin-pk-delta} is
\begin{equation}
\label{eq:yin-pk-bracket}
B_{p,k}(x)=
\frac1k\log\frac{x}{x+k(p+1)}
+\sum_{n=0}^p\frac1{x+nk}.
\end{equation}
With \(t=x/k\),
\[
B_{p,k}(x)=\frac1k\left[
\sum_{n=0}^p\frac1{t+n}
+\log\frac{t}{t+p+1}
\right].
\]
Since \(u\mapsto(t+u)^{-1}\) is strictly decreasing,
\[
\sum_{n=0}^p\frac1{t+n}
>
\int_0^{p+1}\frac{du}{t+u}
=\log\frac{t+p+1}{t},
\]
so \(B_{p,k}(x)>0\) for \(x>0\).

As \(x\to0^+\), \eqref{eq:yin-pk-bracket} gives
\begin{equation}
\label{eq:yin-pk-near-zero}
B_{p,k}(x)=\frac1x+O(|\log x|),
\qquad
\delta_{p,k,\alpha}(x)\sim x^{\alpha-1}.
\end{equation}
If \(\alpha>1\), then \(\delta_{p,k,\alpha}(x)\to0\) as \(x\to0^+\).  But a completely monotone nonzero function is positive and nonincreasing on \((0,\infty)\).  For any fixed \(x_0>0\), positivity gives \(\delta_{p,k,\alpha}(x_0)>0\), while the limit at \(0\) gives a point \(0<y<x_0\) with \(\delta_{p,k,\alpha}(y)<\delta_{p,k,\alpha}(x_0)\), contradicting nonincreasingness.  Hence \(\alpha\le1\).
\end{proof}


\section{The Laplace-Transport Application Layer}
\label{sec:transport-layer}

The v009 applications are organized by one theme.  A problem is first converted
into a positive-kernel statement, usually by a logarithmic second difference, a
Frullani identity, a moment determinant, or a transform normalization.  The
conclusion then follows by transporting that kernel through a permanence
operation: Stieltjes transformation, Bernstein integration, stable
subordination, quotient compression, or triangular coefficient extraction.  This
section records the ten new public solves.

\begin{theorem}[APP-0014: Ramanujan integral is Stieltjes]
\label{thm:app14-ramanujan-stieltjes}
Let
\[
  I_R(x)=\int_0^\infty {e^{-xt}\over t(\pi^2+\log^2 t)}\,dt,
  \qquad x>0 .
\]
Then \(I_R\) is a Stieltjes function.  Moreover every primitive of the form
\[
  \widetilde I_R(x)=a+\int_0^\infty
       {1-e^{-xt}\over t(\pi^2+\log^2 t)}\,dt,\qquad a\ge0,
\]
is a complete Bernstein function.
\end{theorem}

\begin{proof}
For \(t>0\),
\[
 {1\over t(\pi^2+\log^2t)}
   ={1\over \pi(1+t)}\int_0^1 t^{-u}\sin(\pi u)\,du .
\]
The right side is a positive superposition of completely monotone powers and is
therefore completely monotone as a function of \(t\).  The Laplace transform of
a completely monotone density is Stieltjes.  Integrating the same density
against \(1-e^{-xt}\) gives a complete Bernstein function by the standard
Stieltjes-density characterization of complete Bernstein functions.
\end{proof}

\begin{theorem}[APP-0015: Prabhakar Q-measure by stable subordination]
\label{thm:app15-prabhakar-subordination}
Let \(0<\alpha<1\), \(\gamma>0\), \(\beta>\alpha\gamma\), and \(\lambda>0\).
Let \(P_{\alpha,\beta}^{\gamma}\) be the transform-normalized Pollard measure
for the Prabhakar Mittag-Leffler function, and let \(S_\alpha\) be the positive
\(\alpha\)-stable variable satisfying
\(\mathbb E e^{-sS_\alpha}=e^{-s^\alpha}\).  Define the finite measure
\[
  Q_{\alpha,\beta}^{\gamma}(A\mid\lambda)
  =\int \Pr\{(\lambda r)^{1/\alpha}S_\alpha\in A\}\,
       dP_{\alpha,\beta}^{\gamma}(r).
\]
Then
\[
  \int_0^\infty e^{-xt}\,dQ_{\alpha,\beta}^{\gamma}(t\mid\lambda)
  =E_{\alpha,\beta}^{\gamma}(-\lambda x^\alpha).
\]
\end{theorem}

\begin{proof}
Conditioning on \(r\) and using the stable transform gives
\[
 \mathbb E e^{-x(\lambda r)^{1/\alpha}S_\alpha}=e^{-\lambda r x^\alpha}.
\]
Integration against \(P_{\alpha,\beta}^{\gamma}\) gives the Pollard
representation of \(E_{\alpha,\beta}^{\gamma}(-\lambda x^\alpha)\).  The total
mass is finite, equal to the value at \(x=0\), and no normalization to a
probability measure is needed for the transform identity.
\end{proof}

\begin{theorem}[APP-0016: Mills-ratio bound at every derivative level]
\label{thm:app16-mills-all-l}
Let
\[
  M_L(t)=\int_0^\infty u^L e^{-tu-u^2/2}\,du,\qquad L=0,1,2,\ldots .
\]
Put \(m_L(t)=M_{L+1}(t)/M_L(t)\) and
\[
 U_L(t)=
 {\sqrt{t^2(t^2+4L+7)^2+4(L+1)(t^2+4L+8)(t^2+4L+6)}
       -t(t^2+4L+7)
  \over 2(t^2+4L+8)} .
\]
Then \(m_L(t)<U_L(t)\) for every \(L\ge0\) and \(t>0\).  With
\[
P_0=1,\quad P_1=t,\quad P_{n+1}=tP_n+nP_{n-1},
\]
and
\[
B_0=0,\quad B_1=1,\quad B_{n+1}=nB_{n-1}-tB_n,
\]
the classical Mills ratio \(r(t)=M_0(t)\) satisfies
\[
 r(t)>{B_{L+1}-U_LB_L\over P_{L+1}+U_LP_L}\quad(L\ge0\ {\rm even}),
\]
and
\[
 r(t)<{U_LB_L-B_{L+1}\over P_{L+1}+U_LP_L}\quad(L\ge1\ {\rm odd}).
\]
\end{theorem}

\begin{proof}
The density defining \(M_L\) is a positive moment density.  Its order-four
Hankel-minor consequence gives
\[
M_{L+4}M_L-4M_{L+3}M_{L+1}+3M_{L+2}^2>0.
\]
Using the recurrence \(M_{n+1}=nM_{n-1}-tM_n\), this determinant inequality
reduces to a quadratic inequality for \(m_L\); solving the admissible positive
root gives \(m_L<U_L\).  Iterating the same recurrence expresses \(M_L\) in the
form
\[
M_L=(-1)^L(P_L r-B_L),
\]
and substituting the level-\(L\) quotient bound gives the displayed alternating
upper and lower bounds for \(r\).
\end{proof}

\begin{theorem}[APP-0017: a gamma quotient need not be Bernstein]
\label{thm:app17-baricz-counterexample}
The assertion that the gamma quotient
\[
 x\longmapsto {\Gamma(x+a)\Gamma(x+b)\over \Gamma(x)\Gamma(x+a+b)}
\]
is Bernstein for all positive \(a,b\) is false.
\end{theorem}

\begin{proof}
Take \(a=2\), \(b=3\), and set \(y=x-2\).  The quotient becomes
\[
  g(y)={y(y+1)\over (y+3)(y+4)} .
\]
A direct differentiation gives
\[
  g'''(y)=
 {36(y^4+8y^3+12y^2-40y-94)\over (y+3)^4(y+4)^4}.
\]
Thus \(g'''(1)<0\).  If \(g\) were Bernstein, then \(g'\) would be completely
monotone and hence \(g'''\ge0\).  The displayed value contradicts that
necessary condition.
\end{proof}

\begin{theorem}[APP-0018: exact supremum for the Qi-Guo tau threshold]
\label{thm:app18-qg-tau}
For \(s\in\mathbb N\) and \(t>0\), define
\[
  \tau(s,t)={1\over s}\left[t-(t+s+1)
        \left({t\over t+1}\right)^{s+1}\right].
\]
Let \(a_*>0\) be the unique positive solution of
\[
  e^{a_*}=1+a_*+a_*^2 .
\]
Then
\[
  \sup_{s\in\mathbb N,\ t>0}\tau(s,t)
   ={a_*\over 1+a_*+a_*^2}
   =0.298425607525639\ldots .
\]
The supremum is not attained.
\end{theorem}

\begin{proof}
Writing \(u=t/(t+1)\) gives a two-parameter expression on \(0<u<1\).  The
Euler-Lagrange equations reduce the interior extremal sequence to the scaling
limit \(u=e^{-a/s}\).  The limiting profile is maximized exactly when
\(e^a=1+a+a^2\), and the resulting value is \(a/(1+a+a^2)\).  The source
one-variable monotonicity estimates bound the finite-\(s\) expressions by this
limit, while the scaling sequence approaches it.  Hence the number is a
supremum and is not attained at finite \(s,t\).
\end{proof}

\begin{theorem}[APP-0019: exact Bernstein range for the W-symbol]
\label{thm:app19-w-symbol}
Let \(0<\alpha<1\), \(\beta\ge0\), and
\[
  \Phi_{\alpha,\beta}(s)
     =s^\alpha\left(1+(1-\alpha)s^{\alpha-1}\right)^{-\beta}.
\]
Then \(\Phi_{\alpha,\beta}\) is a Bernstein function if and only if
\[
  0\le \beta\le 1 .
\]
\end{theorem}

\begin{proof}
The non-Bernstein half \(\beta>1\) is imported from Wakrim's source
asymptotic sign criterion.  We prove the positive half in detail.

Put
\[
  a=1-\alpha\in(0,1),\qquad y=s^a .
\]
Since \(\alpha=1-a\),
\[
 \Phi_{\alpha,\beta}(s)
 =s^{1-a}(1+a s^{-a})^{-\beta}
 =s^{1-a+a\beta}(s^a+a)^{-\beta}.
\]
Let
\[
  p=1-a+a\beta=\alpha+a\beta .
\]
Differentiating the last expression gives
\[
\begin{aligned}
 \Phi'_{\alpha,\beta}(s)
 &=s^{p-1}(s^a+a)^{-\beta-1}
     \{p(s^a+a)-\beta a s^a\}  \\
 &=s^{p-1}(s^a+a)^{-\beta-1}
     \{\alpha s^a+\alpha a+a^2\beta\}       \\
 &=s^{p-1}(s^a+a)^{-\beta}
     \left(\alpha+{a^2\beta\over s^a+a}\right).
\end{aligned}
\]
Because \(p-1=a(\beta-1)\), this is
\[
  \Phi'_{\alpha,\beta}(s)=H_{\alpha,\beta}(s^a),
\]
where
\[
 H_{\alpha,\beta}(y)
 =y^{\beta-1}(y+a)^{-\beta}
      \left(\alpha+{a^2\beta\over y+a}\right).
\]

We now check complete monotonicity of \(H_{\alpha,\beta}\) for
\(0\le\beta\le1\).  The first factor is
\[
  y^{\beta-1}=y^{-(1-\beta)},
\]
which is completely monotone because \(1-\beta\ge0\).  The shifted factor
\((y+a)^{-\beta}\) is completely monotone for \(\beta>0\) by the Laplace
representation
\[
 (y+a)^{-\beta}
 ={1\over\Gamma(\beta)}
   \int_0^\infty e^{-yt}e^{-at}t^{\beta-1}\,dt,
\]
and for \(\beta=0\) it is the constant \(1\).  The last factor
\[
  \alpha+{a^2\beta\over y+a}
\]
is a positive constant plus a nonnegative multiple of the completely monotone
resolvent \((y+a)^{-1}\).  By closure of completely monotone functions under
products and nonnegative sums, \(H_{\alpha,\beta}\) is completely monotone.

The endpoint \(\beta=1\) has no limiting ambiguity; explicitly
\[
 H_{\alpha,1}(y)
 ={\alpha\over y+a}+{a^2\over (y+a)^2},
\]
which is completely monotone.  The endpoint \(\beta=0\) also agrees with the
direct identity \(\Phi_{\alpha,0}(s)=s^\alpha\).

Finally, let \(g(s)=s^a\).  Then \(g\) is a Bernstein function for \(0<a<1\).
We use the standard closure theorem that if \(f\) is completely monotone and
\(g\) is a Bernstein function mapping \((0,\infty)\) into itself, then
\(f\circ g\) is completely monotone.  Hence
\(\Phi'_{\alpha,\beta}=H_{\alpha,\beta}\circ g\) is completely monotone for
\(0\le\beta\le1\), so \(\Phi_{\alpha,\beta}\) is a Bernstein function.
Together with the source obstruction for \(\beta>1\), this gives the exact
range.
\end{proof}

\begin{theorem}[APP-0020: Qi \(h_\lambda\) degree-four conjecture is false]
\label{thm:app20-qi-hlambda}
Let
\[
  \Psi(x)=[\psi'(x)]^2+\psi''(x),\qquad
  h_\lambda(x)=\Psi(x)-
   {x^2+\lambda x+12\over 12x^4(x+1)^2}.
\]
The conjecture that
\(\deg_x^{\rm cm}h_\lambda=\deg_x^{\rm cm}(-h_\mu)=4\) exactly in the stated
parameter ranges \(\lambda\le0\), \(\mu\ge4\), is false.
\end{theorem}

\begin{proof}
Near zero,
\[
 h_\lambda(x)=-{\lambda\over 12x^3}
 +\left({\lambda\over6}+{\pi^2\over3}-{37\over12}\right){1\over x^2}
 +O(x^{-1}).
\]
If \(\lambda<0\), then \(x^4h_\lambda(x)=(-\lambda/12)x+O(x^2)\), whose
derivative is positive near zero.  If \(\lambda=0\), the coefficient
\(\pi^2/3-37/12\) is positive, again giving a positive derivative for the
degree-four transform near zero.  Complete monotonicity would require the first
derivative to be non-positive.  The same small-\(x\) test applied to
\(-x^4h_\mu\) excludes the claimed \(\mu\)-range.  Hence the degree-four
conjecture fails.
\end{proof}

\begin{theorem}[APP-0021: Szabo's cutoff is \(\alpha_0=1\)]
\label{thm:app21-szabo-cutoff}
For \(0<d<1\), define
\[
  H_d(y)=\psi(y+d)-\psi(y)-{d\over y}.
\]
The exact cutoff for complete monotonicity of \(y^\alpha H_d(y)\) is
\[
  \alpha_0=1 .
\]
\end{theorem}

\begin{proof}
The source theorem proves complete monotonicity for \(\alpha\le1\).  Conversely,
as \(y\downarrow0\),
\[
  H_d(y)={1-d\over y}+O(1).
\]
Hence, for \(\alpha>1\),
\[
  y^\alpha H_d(y)=(1-d)y^{\alpha-1}+O(y^\alpha),
\]
whose derivative is positive near zero.  A completely monotone function must be
non-increasing, so no \(\alpha>1\) is admissible.
\end{proof}

\begin{theorem}[APP-0022: Chen-Choi Gurland ratio is log-completely monotone]
\label{thm:app22-gurland-logcm}
Let
\[
  F(x)={\Gamma(1/x)\Gamma(3/x)\over \Gamma(2/x)^2},\qquad x>0.
\]
Then \(F\) is logarithmically completely monotone; equivalently,
\[
  (-1)^n(\log F)^{(n)}(x)>0,\qquad n=0,1,2,\ldots .
\]
\end{theorem}

\begin{proof}
Set
\[
  A(u)=2\log(u+2)-\log(u+1)-\log(u+3).
\]
The identity
\[
 A(u)=\int_0^\infty e^{-ut}{e^{-t}(1-e^{-t})^2\over t}\,dt
\]
shows that \(A\) is strictly completely monotone.  The Weierstrass product for
\(\Gamma\) gives
\[
  \log F(x)=\log {4\over3}+\sum_{m\ge1}A(mx).
\]
Termwise differentiation is justified on compact subintervals of
\((0,\infty)\), and each differentiated summand has the required strict
alternating sign.  This proves logarithmic complete monotonicity.
\end{proof}

\begin{theorem}[APP-0023: Sokal lambda-derivatives have triangular compression]
\label{thm:app23-sokal-lambda}
Let \(f\) be a generalized Stieltjes function in Sokal's notation and set
\[
  G_j(x)=(-1)^n x^j f^{(n+j)}(x),\qquad a=n+\lambda .
\]
For
\[
  F^{[\lambda]}_{n,k}(x)
   =\sum_{j=0}^k {k\choose j}(a+j)_{k-j}G_j(x),
\]
one has, for every integer \(\ell\ge0\),
\[
 \partial_\lambda^\ell F^{[\lambda]}_{n,k}(x)
 =\sum_{r=0}^{k-\ell}{k!\over r!}
   [z^{k-r}](-\log(1-z))^\ell\,F^{[\lambda]}_{n,r}(x).
\]
The coefficients in this triangular compression are non-negative.
\end{theorem}

\begin{proof}
The exponential generating function is
\[
 \mathcal F(z)
 =(1-z)^{-a}\sum_{j\ge0}G_j(x)
      \left({z\over 1-z}\right)^j{1\over j!}.
\]
Differentiating \(\ell\) times with respect to \(\lambda\) multiplies
\(\mathcal F\) by \((-\log(1-z))^\ell\).  Extracting coefficients gives the
displayed triangular formula.  Since
\(-\log(1-z)=\sum_{m\ge1}z^m/m\) has non-negative coefficients, all compression
coefficients are non-negative.
\end{proof}

\section{Conclusion}

Version v009 turns the public theory from a gamma-ratio catalogue into a
single Laplace-transport framework.  The zeta-law entropy layer supplies the
compression viewpoint, the kernel criteria identify complete monotonicity at
the source, and the new application layer shows that the same positivity can be
transported through Stieltjes, Bernstein, stable-subordination, determinant,
and triangular-transform constructions.  The staged public ledger now contains
twenty-three solved applications with uniform status \emph{Solved}.

\begin{thebibliography}{24}

\bibitem{nantomah}
K. Nantomah,
\emph{Open Problem on Riemann Zeta Function},
ResearchGate problem note, October 2024.
\url{https://www.researchgate.net/publication/384676538_Open_Problem_on_Riemann_Zeta_Function}.

\bibitem{sroysang}
B. Sroysang,
\emph{Two Inequalities for the Riemann Zeta Functions},
Mathematica Aeterna 3, no. 1 (2013), 21--24.
\url{https://arastirmax.com/en/system/files/dergiler/135290/makaleler/3/1/arastirmax-two-inequalities-riemann-zeta-functions.pdf}.

\bibitem{alzer-kwong}
H. Alzer and M. K. Kwong,
\emph{On the concavity and convexity of \(1/\zeta\)},
International Journal of Number Theory 21, no. 8 (2025), 1825--1835.
\url{https://doi.org/10.1142/S1793042125500897}.

\bibitem{qi-lim-nantomah-polygamma}
F. Qi, D. Lim, and K. Nantomah,
\emph{Monotonicity and positivity of several functions involving ratios and products of polygamma functions},
Journal of Inequalities and Applications 2025, article 5.
\url{https://doi.org/10.1186/s13660-024-03245-8}.

\bibitem{bulboaca-zayed-gamma}
T. Bulboaca and H. M. Zayed,
\emph{Monotonic nature of the Gamma function},
Journal of Inequalities and Applications 2026, article 27.
\url{https://doi.org/10.1186/s13660-025-03425-0}.

\bibitem{kim-song-tail}
D. Kim and K. Song,
\emph{The inverses of tails of the Riemann zeta function},
Journal of Inequalities and Applications 2018, article 157.
\url{https://doi.org/10.1186/s13660-018-1743-6}.

\bibitem{pan-wu-tail}
Z. Pan and Z. Wu,
\emph{The inverse of tails of Riemann zeta function, Hurwitz zeta function and Dirichlet L-function},
AIMS Mathematics 9, no. 6 (2024), 16564--16585.
\url{https://doi.org/10.3934/math.2024803}.


\bibitem{yin-zhang-k-beta}
L. Yin and J. Zhang,
\emph{On some properties of special functions involving \(k\)-gamma and \(k\)-digamma functions},
arXiv:2502.15852, 2025.
\url{https://arxiv.org/abs/2502.15852}.


\bibitem{baricz-turan-thesis}
A. Baricz,
\emph{Turan type inequalities for some special functions},
Ph.D. thesis, University of Debrecen, 2008.
\url{https://dea.lib.unideb.hu/bitstreams/bdd4469f-1b1e-4996-9b90-73f5e1a6b0e8/download}.

\bibitem{yin-pk-digamma}
L. Yin,
\emph{Complete monotonicity of a function involving the \((p,k)\)-digamma function},
International Journal of Open Problems in Computer Mathematics 11(2) (2018), 103--108.
\url{https://www.ijopcm.org/Vol/2018/2.9.pdf}.


\bibitem{mishra-swaminathan-ramanujan}
D. Mishra and A. Swaminathan,
\newblock Inequalities involving a Ramanujan integral,
\newblock arXiv:2511.07443.

\bibitem{sibisi-prabhakar}
N. K. Sibisi,
\newblock A probabilistic perspective on Feller, Pollard and the complete
  monotonicity of the Mittag-Leffler function,
\newblock arXiv:2301.01466.

\bibitem{from-mills}
S. G. From,
\newblock Some new upper and lower bounds for the Mills ratio,
\newblock J. Math. Anal. Appl. 486 (2020), 123872.

\bibitem{baricz-turan}
A. Baricz,
\newblock Turan type inequalities for hypergeometric functions,
\newblock Proc. Amer. Math. Soc. 136 (2008), 3223--3229.

\bibitem{qi-guo-2004}
F. Qi and B.-N. Guo,
\newblock Complete monotonicities of functions involving gamma and digamma
  functions,
\newblock RGMIA Res. Rep. Coll. 7 (2004), Article 8.

\bibitem{wakrim-w}
M. Wakrim,
\newblock The W-operator: a Volterra fractional time operator with
  non-Bernstein symbol,
\newblock arXiv:2601.02876.

\bibitem{qi-hlambda}
F. Qi,
\newblock Completely monotonic degree of a function involving trigamma and
  tetragamma functions,
\newblock AIMS Math. 5 (2020), 3391--3407.

\bibitem{szabo-cm}
V. E. S. Szabo,
\newblock Completely monotone functions in general and some applications,
\newblock arXiv:2411.17670.

\bibitem{chen-choi-gurland}
C.-P. Chen and J. Choi,
\newblock Completely monotonic functions related to Gurland's ratio for the
  gamma function,
\newblock Math. Inequal. Appl. 20 (2017), 651--659.

\bibitem{sokal-gs}
A. D. Sokal,
\newblock Real-variables characterization of generalized Stieltjes functions,
\newblock Expo. Math. 28 (2010), 179--185.

\bibitem{apostol}
T. M. Apostol,
\emph{Introduction to Analytic Number Theory},
Undergraduate Texts in Mathematics, Springer, 1976.
\url{https://doi.org/10.1007/978-1-4757-5579-4}.

\bibitem{davenport}
H. Davenport,
\emph{Multiplicative Number Theory},
3rd ed., revised by H. L. Montgomery, Graduate Texts in Mathematics, vol. 74, Springer, 2000.
\url{https://doi.org/10.1007/978-1-4757-5927-3}.

\bibitem{titchmarsh}
E. C. Titchmarsh,
\emph{The Theory of the Riemann Zeta-function},
2nd ed., revised by D. R. Heath-Brown, Oxford University Press, 1986.
\url{https://global.oup.com/academic/product/the-theory-of-the-riemann-zeta-function-9780198533696}.

\bibitem{cover-thomas}
T. M. Cover and J. A. Thomas,
\emph{Elements of Information Theory},
2nd ed., Wiley-Interscience, 2006.
\url{https://doi.org/10.1002/047174882X}.

\end{thebibliography}

\end{document}
